\input{preamble}

% OK, start here.
%
\begin{document}

\title{Espaces algébriques}


\maketitle

\phantomsection
\label{section-phantom}

\tableofcontents

\section{Introduction}
\label{section-introduction}

\noindent
Les espaces algébriques ont été introduits par Michael Artin,
voir \cite{ArtinI}, \cite{ArtinII},
\cite{Artin-Theorem-Representability},
\cite{Artin-Construction-Techniques},
\cite{Artin-Algebraic-Spaces},
\cite{Artin-Algebraic-Approximation},
\cite{Artin-Implicit-Function},
et \cite{ArtinVersal}.
Une partie des fondements a été élaborée conjointement avec
Knutson, qui en a tiré l'ouvrage \cite{Kn}.
Artin a défini (voir \cite[Definition 1.3]{Artin-Implicit-Function})
un espace algébrique comme un faisceau pour la topologie étale
qui est localement représentable pour la topologie étale.
Dans la plupart des travaux d'Artin, les catégories de schémas
considérées sont formées de schémas localement de type fini sur une base
noethérienne excellente fixée.

\medskip\noindent
Notre définition diffère légèrement de la définition originelle d'Artin.
En effet, nos espaces algébriques sont des faisceaux pour la topologie fppf
dont la diagonale est représentable et qui possèdent un «~revêtement~» étale
par un schéma. Travailler avec la topologie fppf plutôt qu'avec la topologie
étale n'est qu'un point technique et ne change presque
rien; nous montrerons dans
Amorçage, Section \ref{bootstrap-section-spaces-etale}
que nous aurions obtenu la même catégorie d'espaces algébriques
en travaillant avec la topologie étale. Dans ce même chapitre,
nous démontrerons que la condition portant sur la diagonale
peut, en un certain sens, être supprimée; voir
Amorçage, Section \ref{bootstrap-section-bootstrap}.

\medskip\noindent
Après avoir défini les espaces algébriques, nous formulons quelques
observations fondamentales. Le principal résultat de ce chapitre est qu'avec
nos définitions, la donnée d'un espace algébrique équivaut à celle d'une
relation d'équivalence étale; voir l'exposé de la Section
\ref{section-presentations} et le Théorème \ref{theorem-presentation}.
L'analogue de ce théorème dans le cadre d'Artin est \cite[Theorem 1.5]{Artin-Implicit-Function}, ou
\cite[Proposition II.1.7]{Kn}. Autrement dit, le faisceau
défini par une relation d'équivalence étale a une diagonale représentable.
Il s'ensuit que notre définition coïncide, au sens large, avec la définition
originelle d'Artin. Cela signifie aussi que l'on peut donner des exemples
d'espaces algébriques en écrivant simplement une relation d'équivalence étale.

\medskip\noindent
Dans la Section \ref{section-separation}, nous introduisons divers axiomes
de séparation pour les espaces algébriques rencontrés dans la littérature.
Enfin, dans la Section \ref{section-examples},
nous donnons des exemples insolites d'espaces algébriques et d'autres qui le sont moins.






\section{Remarques générales}
\label{section-general}

\noindent
Nous travaillons dans un grand site fppf convenable $\Sch_{fppf}$
comme dans Topologies, Définition \ref{topologies-definition-big-fppf-site}.
Ainsi, sauf mention expresse du contraire, tous les schémas seront des objets
de $\Sch_{fppf}$.
Dans la Section \ref{section-change-big-site}, nous étudions les modifications
entraînées par un changement de grand site fppf.

\medskip\noindent
Nous travaillerons toujours relativement à un schéma de base $S$ appartenant à $\Sch_{fppf}$.
Nous travaillerons alors avec le grand site fppf $(\Sch/S)_{fppf}$,
voir Topologies, Définition \ref{topologies-definition-big-small-fppf}.
On retrouve le cas absolu en prenant
$S = \Spec(\mathbf{Z})$.

\medskip\noindent
Si $U, T$ sont des schémas sur $S$, nous désignons
par $U(T)$ l'ensemble des points à valeurs dans $T$ {\it au-dessus} de $S$.
Autrement dit : $U(T) = \Mor_S(T, U)$.

\medskip\noindent
Notons que tout recouvrement fpqc est un épimorphisme effectif universel, voir
Descente, Lemme \ref{descent-lemma-fpqc-universal-effective-epimorphisms}.
Par conséquent, la topologie sur $\Sch_{fppf}$
est moins fine que la topologie canonique, et tous les préfaisceaux représentables
sont des faisceaux.







\section{Morphismes représentables de préfaisceaux}
\label{section-representable}

\noindent
Soit $S$ un schéma appartenant à $\Sch_{fppf}$.
Soient $F, G : (\Sch/S)_{fppf}^{opp} \to \textit{Sets}$.
Soit $a : F \to G$ un morphisme représentable de foncteurs, voir
Catégories,
Définition \ref{categories-definition-representable-map-presheaves}.
Cela signifie que pour tout
$U \in \Ob((\Sch/S)_{fppf})$
et tout $\xi \in G(U)$, le produit fibré $h_U \times_{\xi, G} F$ est représentable.
Choisissons un objet $V_\xi$ qui le représente et un isomorphisme
$h_{V_\xi} \to h_U \times_G F$.
D'après le lemme de Yoneda, voir Catégories, Lemme \ref{categories-lemma-yoneda},
la projection $h_{V_\xi} \to h_U \times_G F \to h_U$ provient d'un unique
morphisme de schémas $a_\xi : V_\xi \to U$.
Le diagramme suivant permet de se représenter la situation :
$$
\xymatrix{
V_\xi \ar@{~>}[r] \ar[d]_{a_\xi} & h_{V_\xi} \ar[d] \ar[r] & F \ar[d]^a \\
U \ar@{~>}[r] & h_U \ar[r]^\xi & G
}
$$
où les flèches ondulées représentent le plongement de Yoneda.
Voici quelques lemmes sur cette notion, valables dans un cadre très général.

\begin{lemma}
\label{lemma-morphism-schemes-gives-representable-transformation}
Soit $S$ un schéma appartenant à $\Sch_{fppf}$ et soient
$X$, $Y$ des objets de $(\Sch/S)_{fppf}$.
Soit $f : X \to Y$ un morphisme de schémas.
Alors
$$
h_f : h_X \longrightarrow h_Y
$$
est un morphisme représentable de foncteurs.
\end{lemma}

\begin{proof}
Cela est formel et repose seulement sur le fait que
la catégorie $(\Sch/S)_{fppf}$ admet des produits fibrés.
\end{proof}

\begin{lemma}
\label{lemma-composition-representable-transformations}
Soit $S$ un schéma appartenant à $\Sch_{fppf}$.
Soient $F, G, H : (\Sch/S)_{fppf}^{opp} \to \textit{Sets}$.
Soient $a : F \to G$, $b : G \to H$ des morphismes représentables de foncteurs.
Alors
$$
b \circ a : F \longrightarrow H
$$
est un morphisme représentable de foncteurs.
\end{lemma}

\begin{proof}
Cela est entièrement formel et vaut dans toute catégorie.
\end{proof}

\begin{lemma}
\label{lemma-base-change-representable-transformations}
Soit $S$ un schéma appartenant à $\Sch_{fppf}$.
Soient $F, G, H : (\Sch/S)_{fppf}^{opp} \to \textit{Sets}$.
Soit $a : F \to G$ un morphisme représentable de foncteurs.
Soit $b : H \to G$ un morphisme de foncteurs quelconque.
Considérons le diagramme de produit fibré
$$
\xymatrix{
H \times_{b, G, a} F \ar[r]_-{b'} \ar[d]_{a'} & F \ar[d]^a \\
H \ar[r]^b & G
}
$$
Alors le changement de base $a'$ est un morphisme représentable de foncteurs.
\end{lemma}

\begin{proof}
Cela est entièrement formel et vaut dans toute catégorie.
\end{proof}

\begin{lemma}
\label{lemma-product-representable-transformations}
Soit $S$ un schéma appartenant à $\Sch_{fppf}$.
Soient $F_i, G_i : (\Sch/S)_{fppf}^{opp} \to \textit{Sets}$, $i = 1, 2$.
Soient $a_i : F_i \to G_i$, $i = 1, 2$,
des morphismes représentables de foncteurs.
Alors
$$
a_1 \times a_2 : F_1 \times F_2 \longrightarrow G_1 \times G_2
$$
est un morphisme représentable de foncteurs.
\end{lemma}

\begin{proof}
Écrivons $a_1 \times a_2$ comme la composée
$F_1 \times F_2 \to G_1 \times F_2 \to G_1 \times G_2$.
La première flèche est le changement de base de $a_1$ par le morphisme
$G_1 \times F_2 \to G_1$, et la seconde flèche
est le changement de base de $a_2$ par le morphisme
$G_1 \times G_2 \to G_2$. Ce lemme est donc une conséquence
formelle des Lemmes \ref{lemma-composition-representable-transformations}
et \ref{lemma-base-change-representable-transformations}.
\end{proof}

\begin{lemma}
\label{lemma-representable-transformation-to-sheaf}
Soit $S$ un schéma appartenant à $\Sch_{fppf}$.
Soient $F, G : (\Sch/S)_{fppf}^{opp} \to \textit{Sets}$.
Soit $a : F \to G$ un morphisme représentable de foncteurs.
Si $G$ est un faisceau, alors $F$ l'est aussi.
\end{lemma}

\begin{proof}
Soit $\{\varphi_i : T_i \to T\}$ une famille couvrante du site
$(\Sch/S)_{fppf}$.
Soient $s_i \in F(T_i)$ satisfaisant à la condition de recollement.
Alors $\sigma_i = a(s_i) \in G(T_i)$ satisfont elles aussi à la condition de recollement.
Il existe donc un unique $\sigma \in G(T)$ tel
que $\sigma_i = \sigma|_{T_i}$. Par hypothèse,
$F' = h_T \times_{\sigma, G, a} F$ est un préfaisceau représentable
et donc (voir les remarques de la Section \ref{section-general}) un faisceau.
Notons que $(\varphi_i, s_i) \in F'(T_i)$ satisfont
elles aussi à la condition de recollement et proviennent donc d'un unique
$(\text{id}_T, s) \in F'(T)$. Manifestement, $s$ est la section de
$F$ que nous cherchions.
\end{proof}

\begin{lemma}
\label{lemma-representable-transformation-diagonal}
Soit $S$ un schéma appartenant à $\Sch_{fppf}$.
Soient $F, G : (\Sch/S)_{fppf}^{opp} \to \textit{Sets}$.
Soit $a : F \to G$ un morphisme représentable de foncteurs.
Alors $\Delta_{F/G} : F \to F \times_G F$ est représentable.
\end{lemma}

\begin{proof}
Soit $U \in \Ob((\Sch/S)_{fppf})$. Soit
$\xi = (\xi_1, \xi_2) \in (F \times_G F)(U)$.
Posons $\xi' = a(\xi_1) = a(\xi_2) \in G(U)$.
Par hypothèse, il existe un schéma $V$ muni d'un morphisme $V \to U$
qui représente le produit fibré $h_U \times_{\xi', G} F$.
En particulier, les éléments $\xi_1, \xi_2$ définissent des morphismes
$f_1, f_2 : U \to V$ au-dessus de $U$. Puisque $V$ représente le
produit fibré $h_U \times_{\xi', G} F$ et que
$\xi' = a \circ \xi_1 = a \circ \xi_2$
nous voyons que, si $g : U' \to U$ est un morphisme, alors
$$
g^*\xi_1 = g^*\xi_2
\Leftrightarrow
f_1 \circ g = f_2 \circ g.
$$
Autrement dit, nous voyons que $h_U \times_{\xi, F \times_G F} F$
est représenté par $V \times_{\Delta, V \times V, (f_1, f_2)} U$
qui est un schéma.
\end{proof}










\section{Listes de propriétés utiles des morphismes de schémas}
\label{section-lists}

\noindent
Pour faciliter les renvois, les remarques suivantes énumèrent les
propriétés de morphismes qui vérifient certaines des conditions
requises dans des résultats ultérieurs.

\begin{remark}
\label{remark-list-properties-stable-base-change}
Voici une liste de propriétés/types de morphismes
qui sont {\it stables par changement de base quelconque} :
\begin{enumerate}
\item immersions fermées, ouvertes et localement fermées, voir
Schémas, Lemme \ref{schemes-lemma-base-change-immersion},
\item quasi-compact, voir
Schémas, Lemme \ref{schemes-lemma-quasi-compact-preserved-base-change},
\item universellement fermé, voir
Schémas, Définition \ref{schemes-definition-universally-closed},
\item (quasi-)séparé, voir
Schémas, Lemme \ref{schemes-lemma-separated-permanence},
\item monomorphisme, voir
Schémas, Lemme \ref{schemes-lemma-base-change-monomorphism}
\item surjectif, voir
Morphismes, Lemme \ref{morphisms-lemma-base-change-surjective},
\item radiciel, voir
Morphismes, Lemme \ref{morphisms-lemma-universally-injective},
\item affine, voir
Morphismes, Lemme \ref{morphisms-lemma-base-change-affine},
\item quasi-affine, voir
Morphismes, Lemme \ref{morphisms-lemma-base-change-quasi-affine},
\item (localement) de type fini, voir
Morphismes, Lemme \ref{morphisms-lemma-base-change-finite-type},
\item (localement) quasi-fini, voir
Morphismes, Lemme \ref{morphisms-lemma-base-change-quasi-finite},
\item (localement) de présentation finie, voir
Morphismes, Lemme \ref{morphisms-lemma-base-change-finite-presentation},
\item localement de type fini de dimension relative $d$, voir
Morphismes, Lemme \ref{morphisms-lemma-base-change-relative-dimension-d},
\item universellement ouvert, voir
Morphismes, Définition \ref{morphisms-definition-open},
\item plat, voir
Morphismes, Lemme \ref{morphisms-lemma-base-change-flat},
\item syntomique, voir
Morphismes, Lemme \ref{morphisms-lemma-base-change-syntomic},
\item lisse, voir
Morphismes, Lemme \ref{morphisms-lemma-base-change-smooth},
\item non ramifié (resp.\ G-non ramifié), voir
Morphismes, Lemme \ref{morphisms-lemma-base-change-unramified},
\item \'etale, voir
Morphismes, Lemme \ref{morphisms-lemma-base-change-etale},
\item propre, voir
Morphismes, Lemme \ref{morphisms-lemma-base-change-proper},
\item H-projectif, voir
Morphismes, Lemme \ref{morphisms-lemma-H-projective-base-change},
\item (localement) projectif, voir
Morphismes, Lemme \ref{morphisms-lemma-base-change-projective},
\item fini ou entier, voir
Morphismes, Lemme \ref{morphisms-lemma-base-change-finite},
\item fini localement libre, voir
Morphismes, Lemme \ref{morphisms-lemma-base-change-finite-locally-free},
\item universellement submersif, voir
Morphismes, Lemme \ref{morphisms-lemma-base-change-universally-submersive},
\item homéomorphisme universel, voir
Morphismes, Lemme \ref{morphisms-lemma-base-change-universal-homeomorphism}.
\end{enumerate}
Compléter au besoin.
\end{remark}

\begin{remark}
\label{remark-list-properties-stable-composition}
Parmi les propriétés de morphismes qui sont stables par changement de base
(telles qu'elles sont énumérées dans la
Remarque \ref{remark-list-properties-stable-base-change}),
les suivantes sont aussi {\it stables par composition} :
\begin{enumerate}
\item immersions fermées, ouvertes et localement fermées, voir
Schémas, Lemme \ref{schemes-lemma-composition-immersion},
\item quasi-compact, voir
Schémas, Lemme \ref{schemes-lemma-composition-quasi-compact},
\item universellement fermé, voir
Morphismes, Lemme \ref{morphisms-lemma-composition-proper},
\item (quasi-)séparé, voir
Schémas, Lemme \ref{schemes-lemma-separated-permanence},
\item monomorphisme, voir
Schémas, Lemme \ref{schemes-lemma-composition-monomorphism},
\item surjectif, voir
Morphismes, Lemme \ref{morphisms-lemma-composition-surjective},
\item radiciel, voir
Morphismes, Lemme \ref{morphisms-lemma-composition-universally-injective},
\item affine, voir
Morphismes, Lemme \ref{morphisms-lemma-composition-affine},
\item quasi-affine, voir
Morphismes, Lemme \ref{morphisms-lemma-composition-quasi-affine},
\item (localement) de type fini, voir
Morphismes, Lemme \ref{morphisms-lemma-composition-finite-type},
\item (localement) quasi-fini, voir
Morphismes, Lemme \ref{morphisms-lemma-composition-quasi-finite},
\item (localement) de présentation finie, voir
Morphismes, Lemme \ref{morphisms-lemma-composition-finite-presentation},
\item universellement ouvert, voir
Morphismes, Lemme \ref{morphisms-lemma-composition-open},
\item plat, voir
Morphismes, Lemme \ref{morphisms-lemma-composition-flat},
\item syntomique, voir
Morphismes, Lemme \ref{morphisms-lemma-composition-syntomic},
\item lisse, voir
Morphismes, Lemme \ref{morphisms-lemma-composition-smooth},
\item non ramifié (resp.\ G-non ramifié), voir
Morphismes, Lemme \ref{morphisms-lemma-composition-unramified},
\item \'etale, voir
Morphismes, Lemme \ref{morphisms-lemma-composition-etale},
\item propre, voir
Morphismes, Lemme \ref{morphisms-lemma-composition-proper},
\item H-projectif, voir
Morphismes, Lemme \ref{morphisms-lemma-H-projective-composition},
\item fini ou entier, voir
Morphismes, Lemme \ref{morphisms-lemma-composition-finite},
\item fini localement libre, voir
Morphismes, Lemme \ref{morphisms-lemma-composition-finite-locally-free},
\item universellement submersif, voir
Morphismes, Lemme \ref{morphisms-lemma-composition-universally-submersive},
\item homéomorphisme universel, voir
Morphismes, Lemme \ref{morphisms-lemma-composition-universal-homeomorphism}.
\end{enumerate}
Compléter au besoin.
\end{remark}

\begin{remark}
\label{remark-list-properties-fpqc-local-base}
Parmi les propriétés mentionnées qui sont stables par changement de base
(énumérées dans la Remarque \ref{remark-list-properties-stable-base-change}),
les suivantes sont aussi {\it locales sur la base pour la topologie fpqc}
(et a fortiori locales sur la base pour la topologie fppf) :
\begin{enumerate}
\item pour les immersions, c'est le cas pour
\begin{enumerate}
\item les immersions fermées, voir
Descente, Lemme \ref{descent-lemma-descending-property-closed-immersion},
\item les immersions ouvertes, voir
Descente, Lemme \ref{descent-lemma-descending-property-open-immersion}, et
\item les immersions quasi-compactes, voir
Descente,
Lemme \ref{descent-lemma-descending-property-quasi-compact-immersion},
\end{enumerate}
\item quasi-compact, voir
Descente, Lemme \ref{descent-lemma-descending-property-quasi-compact},
\item universellement fermé, voir
Descente, Lemme
\ref{descent-lemma-descending-property-universally-closed},
\item (quasi-)séparé, voir
Descente, Lemmes
\ref{descent-lemma-descending-property-quasi-separated}, et
\ref{descent-lemma-descending-property-separated},
\item monomorphisme, voir
Descente, Lemme \ref{descent-lemma-descending-property-monomorphism},
\item surjectif, voir
Descente, Lemme \ref{descent-lemma-descending-property-surjective},
\item radiciel, voir
Descente, Lemme \ref{descent-lemma-descending-property-universally-injective},
\item affine, voir
Descente, Lemme \ref{descent-lemma-descending-property-affine},
\item quasi-affine, voir
Descente, Lemme \ref{descent-lemma-descending-property-quasi-affine},
\item (localement) de type fini, voir
Descente,
Lemmes \ref{descent-lemma-descending-property-locally-finite-type}, et
\ref{descent-lemma-descending-property-finite-type},
\item (localement) quasi-fini, voir
Descente, Lemme \ref{descent-lemma-descending-property-quasi-finite},
\item (localement) de présentation finie, voir
Descente, Lemmes
\ref{descent-lemma-descending-property-locally-finite-presentation}, et
\ref{descent-lemma-descending-property-finite-presentation},
\item localement de type fini de dimension relative $d$, voir
Descente,
Lemme \ref{descent-lemma-descending-property-relative-dimension-d},
\item universellement ouvert, voir
Descente, Lemme \ref{descent-lemma-descending-property-universally-open},
\item plat, voir
Descente, Lemme \ref{descent-lemma-descending-property-flat},
\item syntomique, voir
Descente, Lemme \ref{descent-lemma-descending-property-syntomic},
\item lisse, voir
Descente, Lemme \ref{descent-lemma-descending-property-smooth},
\item non ramifié (resp.\ G-non ramifié), voir
Descente, Lemme \ref{descent-lemma-descending-property-unramified},
\item \'etale, voir
Descente, Lemme \ref{descent-lemma-descending-property-etale},
\item propre, voir
Descente, Lemme \ref{descent-lemma-descending-property-proper},
\item fini ou entier, voir
Descente, Lemme \ref{descent-lemma-descending-property-finite},
\item fini localement libre, voir
Descente, Lemme \ref{descent-lemma-descending-property-finite-locally-free},
\item universellement submersif, voir
Descente, Lemme \ref{descent-lemma-descending-property-universally-submersive},
\item homéomorphisme universel, voir
Descente, Lemme \ref{descent-lemma-descending-property-universal-homeomorphism}.
\end{enumerate}
Notons que la propriété d'être une ``immersion'' n'est pas nécessairement locale
sur la base pour la topologie fpqc, mais que dans
Descente, Lemme \ref{descent-lemma-descending-fppf-property-immersion},
nous avons démontré qu'elle est locale sur la base pour la topologie fppf.
\end{remark}








\section{Propriétés des morphismes représentables de préfaisceaux}
\label{section-representable-properties}

\noindent
Voici la définition qui permet de procéder.

\begin{definition}
\label{definition-relative-representable-property}
On se donne $S$ et $a : F \to G$ représentable comme ci-dessus.
Soit $\mathcal{P}$ une propriété de morphismes de schémas qui
\begin{enumerate}
\item est stable par tout changement de base,
voir Schémas, Définition \ref{schemes-definition-preserved-by-base-change},
et
\item est locale sur la base pour la topologie fppf, voir
Descente, Définition \ref{descent-definition-property-morphisms-local}.
\end{enumerate}
Dans ce cas, nous disons que $a$ possède {\it la propriété $\mathcal{P}$} si, pour tout
$U \in \Ob((\Sch/S)_{fppf})$ et
tout $\xi \in G(U)$, le morphisme de schémas qui en résulte
$V_\xi \to U$ possède la propriété $\mathcal{P}$.
\end{definition}

\noindent
Il importe de noter que nous n'utiliserons cette définition que pour
des propriétés de morphismes stables par changement de base et
locales sur la base pour la topologie fppf. Ce choix ne vient
pas de ce que la définition n'aurait autrement aucun sens ; il vient plutôt
de ce que nous pouvons vouloir donner une autre définition
mieux adaptée à la propriété envisagée.

\begin{remark}
\label{remark-warning}
Considérons la propriété $\mathcal{P}=$``surjectif''.
Dans ce cas, dire
``soit $F \to G$ un morphisme surjectif'' pourrait être ambigu.
En effet, nous pourrions entendre la notion définie
dans la Définition \ref{definition-relative-representable-property}
ci-dessus, ou bien un morphisme surjectif de préfaisceaux, voir
Sites, Définition \ref{sites-definition-presheaves-injective-surjective},
ou encore, si $F$ et $G$ sont tous deux des faisceaux,
un morphisme surjectif de faisceaux, voir
Sites, Définition \ref{sites-definition-sheaves-injective-surjective}.
Sauf mention expresse du contraire, lorsque nous parlerons de morphismes d'espaces algébriques,
nous entendrons toujours la première notion. Voir le
Lemme \ref{lemma-surjective-flat-locally-finite-presentation}
pour un cas où cette propriété entraîne la surjectivité du morphisme de faisceaux sous-jacent.
\end{remark}

\noindent
Voici une vérification de cohérence.

\begin{lemma}
\label{lemma-morphism-schemes-gives-representable-transformation-property}
Soient $S$, $X$, $Y$ des objets de $\Sch_{fppf}$.
Soit $f : X \to Y$ un morphisme de schémas.
Soit $\mathcal{P}$ comme dans la
Définition \ref{definition-relative-representable-property}.
Alors $h_X \longrightarrow h_Y$ possède la propriété $\mathcal{P}$ si
et seulement si $f$ possède la propriété $\mathcal{P}$.
\end{lemma}

\begin{proof}
Notons que le lemme a un sens en vertu du
Lemme \ref{lemma-morphism-schemes-gives-representable-transformation}.
Démonstration omise.
\end{proof}

\begin{lemma}
\label{lemma-composition-representable-transformations-property}
Soit $S$ un schéma appartenant à $\Sch_{fppf}$.
Soient $F, G, H : (\Sch/S)_{fppf}^{opp} \to \textit{Sets}$.
Soit $\mathcal{P}$ une propriété comme dans la
Définition \ref{definition-relative-representable-property}
qui est stable par composition.
Soient $a : F \to G$, $b : G \to H$ des morphismes représentables de foncteurs.
Si $a$ et $b$ possèdent la propriété $\mathcal{P}$, il en va de même de
$b \circ a : F \longrightarrow H$.
\end{lemma}

\begin{proof}
Notons que le lemme a un sens en vertu du
Lemme \ref{lemma-composition-representable-transformations}.
Démonstration omise.
\end{proof}

\begin{lemma}
\label{lemma-base-change-representable-transformations-property}
Soit $S$ un schéma appartenant à $\Sch_{fppf}$.
Soient $F, G, H : (\Sch/S)_{fppf}^{opp} \to \textit{Sets}$.
Soit $\mathcal{P}$ une propriété comme dans la
Définition \ref{definition-relative-representable-property}.
Soit $a : F \to G$ un morphisme représentable de foncteurs.
Soit $b : H \to G$ un morphisme quelconque de foncteurs.
Considérons le diagramme de produit fibré
$$
\xymatrix{
H \times_{b, G, a} F \ar[r]_-{b'} \ar[d]_{a'} & F \ar[d]^a \\
H \ar[r]^b & G
}
$$
Si $a$ possède la propriété $\mathcal{P}$, alors son changement de base $a'$
possède également la propriété $\mathcal{P}$.
\end{lemma}

\begin{proof}
Notons que le lemme a un sens en vertu du
Lemme \ref{lemma-base-change-representable-transformations}.
Démonstration omise.
\end{proof}

\begin{lemma}
\label{lemma-descent-representable-transformations-property}
Soit $S$ un schéma appartenant à $\Sch_{fppf}$.
Soient $F, G, H : (\Sch/S)_{fppf}^{opp} \to \textit{Sets}$.
Soit $\mathcal{P}$ une propriété comme dans la
Définition \ref{definition-relative-representable-property}.
Soit $a : F \to G$ un morphisme représentable de foncteurs.
Soit $b : H \to G$ un morphisme quelconque de foncteurs.
Considérons le diagramme de produit fibré
$$
\xymatrix{
H \times_{b, G, a} F \ar[r]_-{b'} \ar[d]_{a'} & F \ar[d]^a \\
H \ar[r]^b & G
}
$$
Supposons que $b$ induise un morphisme surjectif de faisceaux pour la topologie fppf $H^\# \to G^\#$.
Dans ce cas, si $a'$ possède la propriété $\mathcal{P}$, alors $a$
possède également la propriété $\mathcal{P}$.
\end{lemma}

\begin{proof}
Remarquons d'abord que, d'après le
Lemme \ref{lemma-base-change-representable-transformations},
le morphisme $a'$ est représentable.
Soient $U \in \Ob((\Sch/S)_{fppf})$ et
$\xi \in G(U)$. Par hypothèse, il existe un recouvrement fppf
$\{U_i \to U\}_{i \in I}$ et des éléments $\xi_i \in H(U_i)$
qui sont envoyés sur $\xi|_{U_i}$ par $b$. Il résulte de la théorie générale des catégories que, pour
chaque $i$, on a un diagramme de produit fibré
$$
\xymatrix{
U_i \times_{\xi_i, H, a'} (H \times_{b, G, a} F) \ar[r] \ar[d] &
U \times_{\xi, G, a} F \ar[d] \\
U_i \ar[r] & U
}
$$
Par hypothèse, la flèche verticale de gauche est un morphisme de schémas qui
possède la propriété $\mathcal{P}$. Comme $\mathcal{P}$ est locale sur la base pour la
topologie fppf, il en résulte que la flèche verticale de droite possède aussi la propriété
$\mathcal{P}$, comme souhaité.
\end{proof}

\begin{lemma}
\label{lemma-product-representable-transformations-property}
Soit $S$ un schéma appartenant à $\Sch_{fppf}$.
Soient $F_i, G_i : (\Sch/S)_{fppf}^{opp} \to \textit{Sets}$,
$i = 1, 2$.
Soient $a_i : F_i \to G_i$, $i = 1, 2$, des morphismes représentables
de foncteurs.
Soit $\mathcal{P}$ une propriété comme dans la
Définition \ref{definition-relative-representable-property}
qui est stable par composition.
Si $a_1$ et $a_2$ possèdent la propriété $\mathcal{P}$, il en va de même de
$a_1 \times a_2 : F_1 \times F_2 \longrightarrow G_1 \times G_2$.
\end{lemma}

\begin{proof}
Notons que le lemme a un sens en vertu du
Lemme \ref{lemma-product-representable-transformations}.
Démonstration omise.
\end{proof}

\begin{lemma}
\label{lemma-representable-transformations-property-implication}
Soit $S$ un schéma appartenant à $\Sch_{fppf}$.
Soient $F, G : (\Sch/S)_{fppf}^{opp} \to \textit{Sets}$.
Soit $a : F \to G$ un morphisme représentable de foncteurs.
Soient $\mathcal{P}$, $\mathcal{P}'$ des propriétés comme dans la
Définition \ref{definition-relative-representable-property}.
Supposons que, pour tout morphisme de schémas $f : X \to Y$,
on ait $\mathcal{P}(f) \Rightarrow \mathcal{P}'(f)$.
Si $a$ possède la propriété $\mathcal{P}$, alors
$a$ possède la propriété $\mathcal{P}'$.
\end{lemma}

\begin{proof}
Formel.
\end{proof}

\begin{lemma}
\label{lemma-surjective-flat-locally-finite-presentation}
Soit $S$ un schéma.
Soient $F, G : (\Sch/S)_{fppf}^{opp} \to \textit{Sets}$ des faisceaux.
Soit $a : F \to G$ représentable, plat,
localement de présentation finie et surjectif.
Alors $a : F \to G$ est surjectif en tant que morphisme de faisceaux.
\end{lemma}

\begin{proof}
Soient $T$ un schéma au-dessus de $S$ et $g : T \to G$ un point à valeurs dans $T$ du foncteur
$G$. Par hypothèse, $T' = F \times_G T$ est (représenté par) un schéma et
le morphisme $T' \to T$ est plat, localement de présentation finie et
surjectif. Ainsi, $\{T' \to T\}$ est un recouvrement fppf tel
que $g|_{T'} \in G(T')$ provient d'un élément de $F(T')$, à savoir
le morphisme $T' \to F$. Cela prouve que le morphisme est surjectif en tant que
morphisme de faisceaux, voir
Sites, Définition \ref{sites-definition-sheaves-injective-surjective}.
\end{proof}

\noindent
Voici une caractérisation des foncteurs dont la
diagonale est représentable.

\begin{lemma}
\label{lemma-representable-diagonal}
Soit $S$ un schéma appartenant à $\Sch_{fppf}$.
Soit $F$ un préfaisceau d'ensembles sur $(\Sch/S)_{fppf}$.
Les assertions suivantes sont équivalentes :
\begin{enumerate}
\item la diagonale $F \to F \times F$ est représentable,
\item pour $U \in \Ob((\Sch/S)_{fppf})$ et tout $a \in F(U)$,
le morphisme $a : h_U \to F$ est représentable,
\item pour toute paire $U, V \in \Ob((\Sch/S)_{fppf})$
et tous $a \in F(U)$ et $b \in F(V)$, le produit fibré
$h_U \times_{a, F, b} h_V$ est représentable.
\end{enumerate}
\end{lemma}

\begin{proof}
C'est purement formel, voir
Catégories, Lemme \ref{categories-lemma-representable-diagonal}.
Cela ne dépend que du fait que la catégorie $(\Sch/S)_{fppf}$
possède les produits de deux objets et les produits fibrés, voir
Topologies, Lemme \ref{topologies-lemma-fibre-products-fppf}.
\end{proof}

\noindent
Dans la situation du lemme, pour tout morphisme
$\xi : h_U \to F$ comme dans son énoncé, il est légitime
de dire que $\xi$ possède la propriété $\mathcal{P}$, pour toute propriété
comme dans la Définition \ref{definition-relative-representable-property}.
C'est notamment le cas pour $\mathcal{P} = $ ``surjectif''
et $\mathcal{P} = $ ``\'etale'', voir la
Remarque \ref{remark-list-properties-fpqc-local-base}
ci-dessus. Nous utiliserons cette remarque dans la définition
des espaces algébriques ci-dessous.

\begin{lemma}
\label{lemma-transformation-diagonal-properties}
Soit $S$ un schéma appartenant à $\Sch_{fppf}$.
Soit $F$ un préfaisceau d'ensembles sur $(\Sch/S)_{fppf}$.
Soit $\mathcal{P}$ une propriété comme dans la
Définition \ref{definition-relative-representable-property}.
Supposons que, pour tous $U, V \in \Ob((\Sch/S)_{fppf})$, $a \in F(U)$ et
$b \in F(V)$, on ait
\begin{enumerate}
\item $h_U \times_{a, F, b} h_V$ est représentable, disons par le schéma $W$, et
\item le morphisme $W \to U \times_S V$ correspondant au morphisme
$h_U \times_{a, F, b} h_V \to h_U \times h_V$ possède la propriété $\mathcal{P}$,
\end{enumerate}
alors $\Delta : F \to F \times F$ est représentable et possède
la propriété $\mathcal{P}$.
\end{lemma}

\begin{proof}
Observons que $\Delta$ est représentable d'après le
Lemme \ref{lemma-representable-diagonal}.
Nous pouvons reformuler la condition (2) en disant que
le morphisme $h_U \times_{a, F, b} h_V \to h_{U \times_S V}$
possède la propriété $\mathcal{P}$, voir le Lemme
\ref{lemma-morphism-schemes-gives-representable-transformation-property}.
Considérons $T \in \Ob((\Sch/S)_{fppf})$ et $(a, b) \in (F \times F)(T)$.
Observons que nous avons le diagramme commutatif
$$
\xymatrix{
F \times_{\Delta, F \times F, (a, b)} h_T \ar[d] \ar[r] &
h_T \ar[d]^{\Delta_{T/S}} \\
h_T \times_{a, F, b} h_T \ar[r] \ar[d] &
h_{T \times_S T} \ar[d]^{(a, b)} \\
F \ar[r]^\Delta & F \times F
}
$$
dont les deux carrés sont cartésiens. Nous voyons ainsi que
le morphisme $F \times_{F \times F} h_T \to h_T$ est obtenu par changement
de base d'un morphisme qui possède la propriété $\mathcal{P}$ au moyen de
$\Delta_{T/S}$. Comme $\mathcal{P}$ est stable par changement de
base, cela achève la démonstration.
\end{proof}


















\section{Espaces algébriques}
\label{section-algebraic-spaces}

\noindent
Voici la définition.

\begin{definition}
\label{definition-algebraic-space}
Soit $S$ un schéma appartenant à $\Sch_{fppf}$.
Un {\it espace algébrique au-dessus de $S$} est un préfaisceau
$$
F : (\Sch/S)^{opp}_{fppf} \longrightarrow \textit{Sets}
$$
qui possède les propriétés suivantes :
\begin{enumerate}
\item Le préfaisceau $F$ est un faisceau.
\item Le morphisme diagonal $F  \to F \times F$ est représentable.
\item Il existe un schéma $U \in \Ob((\Sch/S)_{fppf})$
et un morphisme $h_U \to F$ qui est surjectif et \'etale\footnote{Voir le
Lemme \ref{lemma-representable-diagonal}, la
Définition \ref{definition-relative-representable-property} et la
Remarque \ref{remark-warning}.}.
\end{enumerate}
\end{definition}

\noindent
Notre définition diffère en deux points de la définition ``usuelle'', par exemple de la
définition donnée dans l'ouvrage de Knutson \cite{Kn}.

\medskip\noindent
La première est que nous exigeons que $F$ soit un faisceau pour la topologie fppf.
L'une des raisons de ce choix est que de nombreux exemples naturels
d'espaces algébriques vérifient la condition de faisceau pour les recouvrements fppf
(et même pour les recouvrements fpqc). En outre, l'utilité des espaces
algébriques tient notamment aux résultats de Michael Artin à leur sujet.
Sa méthode comporte une condition qui garantit que le résultat est
localement de présentation finie au-dessus de $S$.
Tout bien considéré, il nous semble que la topologie fppf
est la topologie naturelle à employer. Finalement, on obtient
la même catégorie d'espaces algébriques. Voir
Amorçage, Section \ref{bootstrap-section-spaces-etale}.

\medskip\noindent
La seconde est que nous exigeons seulement que le morphisme diagonal de $F$ soit
représentable, tandis que \cite{Kn} exige qu'il soit également
quasi-compact. Si $F = h_U$ pour un schéma $U$ au-dessus de $S$,
cela revient à demander que $U$ soit quasi-séparé.
Nous cherchons à démontrer un certain
nombre des résultats qui suivent en supposant seulement que la diagonale
de $F$ soit représentable, et à ajouter simplement une hypothèse supplémentaire chaque fois que
cela est nécessaire. Cette convention a en tout cas pour conséquence agréable que
le lemme suivant est vrai.

\begin{lemma}
\label{lemma-scheme-is-space}
Tout schéma est un espace algébrique. Plus précisément,
pour tout schéma $T \in \Ob((\Sch/S)_{fppf})$,
le foncteur représentable $h_T$ est un espace algébrique.
\end{lemma}

\begin{proof}
Le foncteur $h_T$ est un faisceau d'après nos remarques de la Section \ref{section-general}.
La diagonale $h_T \to h_T \times h_T = h_{T \times T}$ est
représentable parce que $(\Sch/S)_{fppf}$ possède les produits fibrés.
Le morphisme identité $h_T \to h_T$ est surjectif et \'etale.
\end{proof}

\begin{definition}
\label{definition-morphism-algebraic-spaces}
Soient $F$, $F'$ des espaces algébriques au-dessus de $S$.
Un {\it morphisme $f : F \to F'$ d'espaces algébriques au-dessus de $S$}
est un morphisme de foncteurs de $F$ vers $F'$.
\end{definition}

\noindent
La catégorie des espaces algébriques au-dessus de $S$ contient la catégorie
$(\Sch/S)_{fppf}$ comme sous-catégorie pleine par l'intermédiaire du
plongement de Yoneda $T/S \mapsto h_T$. Désormais, nous ne distinguerons plus
un schéma $T/S$ de l'espace algébrique qu'il représente.
Ainsi, lorsque nous disons ``Soit $f : T \to F$ un morphisme du schéma
$T$ vers l'espace algébrique $F$'', nous voulons dire que
$T \in \Ob((\Sch/S)_{fppf})$, que $F$ est un
espace algébrique au-dessus de $S$ et que $f : h_T \to F$ est un morphisme
d'espaces algébriques au-dessus de $S$.






\section{Produits fibrés d'espaces algébriques}
\label{section-fibre-products}

\noindent
La catégorie des espaces algébriques au-dessus de $S$ possède à la fois des produits et des
produits fibrés.

\begin{lemma}
\label{lemma-product-spaces}
Soit $S$ un schéma appartenant à $\Sch_{fppf}$.
Soient $F, G$ des espaces algébriques au-dessus de $S$.
Alors $F \times G$ est un espace algébrique et c'est un produit
dans la catégorie des espaces algébriques au-dessus de $S$.
\end{lemma}

\begin{proof}
Il est immédiat que $H = F \times G$ est un faisceau.
La diagonale de $H$ est simplement le produit des
diagonales de $F$ et de $G$. Elle est donc représentable d'après le
Lemme \ref{lemma-product-representable-transformations}.
Enfin, si $U \to F$ et $V \to G$ sont des morphismes étales
surjectifs, où $U, V \in \Ob((\Sch/S)_{fppf})$,
alors $U \times V \to F \times G$ est étale surjectif
d'après le Lemme \ref{lemma-product-representable-transformations-property}.
\end{proof}

\begin{lemma}
\label{lemma-fibre-product-spaces-over-sheaf-with-representable-diagonal}
Soit $S$ un schéma appartenant à $\Sch_{fppf}$.
Soit $H$ un faisceau sur $(\Sch/S)_{fppf}$ dont la diagonale
est représentable. Soient $F, G$ des espaces algébriques au-dessus de $S$.
Soient $F \to H$, $G \to H$ des morphismes de faisceaux.
Alors $F \times_H G$ est un espace algébrique.
\end{lemma}

\begin{proof}
Vérifions les trois conditions de la
Définition \ref{definition-algebraic-space}.
Un produit fibré de faisceaux est un faisceau; ainsi, $F \times_H G$ est un faisceau.
La diagonale de $F \times_H G$ est la flèche verticale de gauche dans le diagramme
$$
\xymatrix{
F \times_H G \ar[r] \ar[d]_\Delta &
F \times G \ar[d]^{\Delta_F \times \Delta_G} \\
(F \times F) \times_{(H \times H)} (G \times G) \ar[r] &
(F \times F) \times (G \times G)
}
$$
qui est cartésien. Ainsi, $\Delta$ est représentable, puisqu'elle s'obtient par changement de base
à partir du morphisme de droite, qui est représentable; voir les
Lemmes \ref{lemma-product-representable-transformations} et
\ref{lemma-base-change-representable-transformations}.
Enfin, soient $U, V \in \Ob((\Sch/S)_{fppf})$
et soient $a : U \to F$, $b : V \to G$ des morphismes étales surjectifs.
Comme $\Delta_H$ est représentable, on voit que $U \times_H V$ est un schéma.
Le morphisme
$$
U \times_H V \longrightarrow F \times_H G
$$
est étale surjectif, car il est le composé des morphismes
$U \times_H V \to U \times_H G$ et $U \times_H G \to F \times_H G$,
obtenus par changement de base à partir des morphismes étales surjectifs $U \to F$ et $V \to G$; voir les
Lemmes \ref{lemma-composition-representable-transformations} et
\ref{lemma-base-change-representable-transformations}.
Cela montre que la dernière condition de la
Définition \ref{definition-algebraic-space}
est satisfaite, et l'on conclut que $F \times_H G$ est un espace algébrique.
\end{proof}

\begin{lemma}
\label{lemma-fibre-product-spaces}
Soit $S$ un schéma appartenant à $\Sch_{fppf}$.
Soient $F \to H$, $G \to H$ des morphismes d'espaces algébriques au-dessus de $S$.
Alors $F \times_H G$ est un espace algébrique et c'est un produit fibré
dans la catégorie des espaces algébriques au-dessus de $S$.
\end{lemma}

\begin{proof}
Il résulte du Lemme plus général
\ref{lemma-fibre-product-spaces-over-sheaf-with-representable-diagonal}
que $F \times_H G$ est un espace algébrique.
Il est clair que $F \times_H G$
est un produit fibré dans la catégorie des espaces algébriques au-dessus de $S$,
puisque celle-ci est une sous-catégorie pleine de la catégorie
des (pré)faisceaux d'ensembles sur $(\Sch/S)_{fppf}$.
\end{proof}






\section{Recollement d'espaces algébriques}
\label{section-glueing-algebraic-spaces}

\noindent
Dans cette section, nous commençons véritablement à abuser des notations et à ne plus
distinguer les schémas des espaces qu'ils représentent.

\begin{lemma}
\label{lemma-coproduct-sheaves-open-and-closed}
Soit $S \in \Ob(\Sch_{fppf})$. Soient $F$ et $G$ des faisceaux sur
$(\Sch/S)_{fppf}^{opp}$ et notons $F \amalg G$ leur somme
dans la catégorie des faisceaux. Le morphisme $F \to F \amalg G$ est représentable par
des immersions à la fois ouvertes et fermées.
\end{lemma}

\begin{proof}
Soient $U$ un schéma et $\xi \in (F \amalg G)(U)$. Rappelons que la
somme dans la catégorie des faisceaux est le faisceau associé au
préfaisceau somme (Sites, Lemme \ref{sites-lemma-colimit-sheaves}).
Il existe donc un recouvrement fppf $\{g_i : U_i \to U\}_{i \in I}$
et une décomposition en somme disjointe $I = I' \amalg I''$ telle que
$U_i \to U \to F \amalg G$ se factorise par $F$, resp.\ par $G$,
si et seulement si $i \in I'$, resp.\ $i \in I''$. Puisque l'intersection de $F$ et de
$G$ dans $F \amalg G$ est vide, on en déduit que
$U_i \times_U U_j$ est vide si $i \in I'$ et $j \in I''$.
Par conséquent, $U' = \bigcup_{i \in I'} g_i(U_i)$ et
$U'' = \bigcup_{i \in I''} g_i(U_i)$ sont des sous-schémas ouverts
disjoints de $U$ (Morphismes, Lemme \ref{morphisms-lemma-fppf-open})
tels que $U = U' \amalg U''$.
Nous omettons la vérification de l'égalité $U' = U \times_{F \amalg G} F$.
\end{proof}

\begin{lemma}
\label{lemma-representable-sheaf-coproduct-sheaves}
Soit $S \in \Ob(\Sch_{fppf})$.
Soit $U \in \Ob((\Sch/S)_{fppf})$.
Étant donnés un ensemble $I$ et des faisceaux $F_i$ sur $\Ob((\Sch/S)_{fppf})$,
si $U \cong \coprod_{i\in I} F_i$
comme faisceaux, chaque $F_i$ est représentable par un sous-schéma
$U_i$ à la fois ouvert et fermé, et $U \cong \coprod U_i$ comme schémas.
\end{lemma}

\begin{proof}
D'après le Lemme \ref{lemma-coproduct-sheaves-open-and-closed},
le morphisme $F_i \to U$ est représentable par des immersions à la fois ouvertes et fermées.
Ainsi, $F_i$ est représentable par un sous-schéma $U_i$ à la fois ouvert et fermé de $U$.
On a $U = \coprod U_i$, puisque $U \cong \coprod F_i$
comme faisceaux et que l'égalité peut être vérifiée sur les points.
\end{proof}

\begin{lemma}
\label{lemma-algebraic-space-coproduct-sheaves}
Soit $S \in \Ob(\Sch_{fppf})$.
Soit $F$ un espace algébrique au-dessus de $S$.
Étant donnés un ensemble $I$ et des faisceaux $F_i$ sur
$\Ob((\Sch/S)_{fppf})$,
si $F \cong \coprod_{i\in I} F_i$ comme faisceaux,
chaque $F_i$ est un espace algébrique au-dessus de $S$.
\end{lemma}

\begin{proof}
La représentabilité de $F \to F \times F$ implique que chacun des morphismes diagonaux
$F_i \to F_i \times F_i$ est représentable (cela résulte immédiatement des définitions
et de l'égalité $F \times_{(F \times F)} (F_i \times F_i) = F_i$).
Choisissons un schéma $U$ dans $(\Sch/S)_{fppf}$ et un morphisme
\'etale surjectif $U \to F$ (dont l'existence résulte de l'hypothèse).
Le morphisme $U \times_F F_i \to F_i$ obtenu par changement de base est \'etale surjectif
d'après le Lemme \ref{lemma-base-change-representable-transformations-property}.
D'autre part, $U \times_F F_i$ est un schéma d'après le
Lemme \ref{lemma-coproduct-sheaves-open-and-closed}.
Nous avons donc vérifié toutes les conditions de la
Définition \ref{definition-algebraic-space},
et $F_i$ est un espace algébrique.
\end{proof}

\noindent
La condition portant sur la taille de $I$ et des $F_i$ dans le
lemme suivant peut être ignorée par ceux que ne préoccupent pas les
questions de théorie des ensembles.

\begin{lemma}
\label{lemma-coproduct-algebraic-spaces}
Soit $S \in \Ob(\Sch_{fppf})$.
Supposons donnés un ensemble $I$ et des espaces algébriques $F_i$, $i \in I$.
Alors $F = \coprod_{i \in I} F_i$ est un espace algébrique,
pourvu que $I$ et les $F_i$ ne soient pas trop «~grands~» : par exemple, si nous
pouvons choisir des morphismes \'etales surjectifs $U_i \to F_i$ tels que
$\coprod_{i \in I} U_i$ soit isomorphe à un objet de
$(\Sch/S)_{fppf}$, alors $F$ est un espace algébrique.
\end{lemma}

\begin{proof}
Par construction, $F$ est un faisceau. Nous omettons de vérifier que le
morphisme diagonal de $F$ est représentable. Enfin, si $U$ est un
objet de $(\Sch/S)_{fppf}$ isomorphe à $\coprod_{i \in I} U_i$,
alors il est immédiat de vérifier que le morphisme qui en résulte
$U \to \coprod F_i$ est \'etale surjectif.
\end{proof}

\noindent
Voici l'analogue de Schémas, Lemme \ref{schemes-lemma-glue-functors}.

\begin{lemma}
\label{lemma-glueing-algebraic-spaces}
Soit $S \in \Ob(\Sch_{fppf})$.
Soit $F$ un préfaisceau d'ensembles sur $(\Sch/S)_{fppf}$.
Supposons que
\begin{enumerate}
\item $F$ est un faisceau,
\item il existe un ensemble d'indices $I$
et des sous-foncteurs $F_i \subset F$ tels que
\begin{enumerate}
\item chaque $F_i$ est un espace algébrique,
\item chaque $F_i \to F$ est représentable,
\item chaque $F_i \to F$ est une immersion ouverte (voir la
Définition \ref{definition-relative-representable-property}),
\item le morphisme $\coprod F_i \to F$ est surjectif comme morphisme de faisceaux, et
\item $\coprod F_i$ est un espace algébrique (condition ensembliste, voir le
Lemme \ref{lemma-coproduct-algebraic-spaces}).
\end{enumerate}
\end{enumerate}
Alors $F$ est un espace algébrique.
\end{lemma}

\begin{proof}
Soit $T$ un objet de $(\Sch/S)_{fppf}$. Soit $T \to F$ un morphisme.
D'après les hypothèses (2)(b) et (2)(c), le produit fibré $F_i \times_F T$
est représentable par un sous-schéma ouvert $V_i \subset T$. Il s'ensuit que
$(\coprod F_i) \times_F T$ est représenté par le schéma $\coprod V_i$ au-dessus de $T$.
D'après l'hypothèse (2)(d), il existe un recouvrement fppf $\{T_j \to T\}_{j \in J}$
tel que $T_j \to T \to F$ se factorise par $F_i$, où $i = i(j)$.
Ainsi, $T_j \to T$ se factorise par le sous-schéma ouvert $V_{i(j)} \subset T$.
Puisque la famille $\{T_j \to T\}$ est surjective, il s'ensuit que
$T = \bigcup V_i$ est un recouvrement ouvert. En particulier, le morphisme
de foncteurs $\coprod F_i \to F$ est représentable
et surjectif au sens de la
Définition \ref{definition-relative-representable-property}
(voir la Remarque \ref{remark-warning} pour une discussion).

\medskip\noindent
Ensuite, soit $T' \to F$ un second morphisme ayant pour source un objet de $(\Sch/S)_{fppf}$.
Écrivons comme ci-dessus $T' = \bigcup V'_i$, où $V'_i = T' \times_F F_i$.
Pour montrer que la diagonale $F \to F \times F$ est représentable,
il faut montrer que $G = T \times_F T'$ est représentable; voir le
Lemme \ref{lemma-representable-diagonal}.
Considérons les sous-foncteurs $G_i = G \times_F F_i$.
Notons que $G_i = V_i \times_{F_i} V'_i$; il est donc représentable,
car $F_i$ est un espace algébrique.
D'après ce qui précède, les $G_i$ forment un recouvrement de Zariski de $G$.
Ainsi, d'après Schémas, Lemme \ref{schemes-lemma-glue-functors},
on voit que $G$ est représentable.

\medskip\noindent
Choisissons un schéma $U \in \Ob((\Sch/S)_{fppf})$ et un morphisme
\'etale surjectif $U \to \coprod F_i$ (dont l'existence résulte de l'hypothèse).
Nous pouvons écrire $U = \coprod U_i$, où $U_i$ est l'image réciproque de $F_i$;
voir le Lemme \ref{lemma-representable-sheaf-coproduct-sheaves}.
Nous affirmons que $U \to F$ est \'etale surjectif. La surjectivité découle
de celle de $\coprod F_i \to F$ (voir le premier paragraphe de la démonstration),
en appliquant le
Lemme \ref{lemma-composition-representable-transformations-property}.
Considérons le produit fibré $U \times_F T$, où $T \to F$ est comme
ci-dessus. Il faut montrer que $U \times_F T \to T$ est \'etale.
Puisque $U \times_F T = \coprod U_i \times_F T$, il suffit de montrer que
chacun des morphismes $U_i \times_F T \to T$ est \'etale. Comme
$U_i \times_F T = U_i \times_{F_i} V_i$, cela résulte du
fait que $U_i \to F_i$ est \'etale et que $V_i \to T$ est une immersion ouverte
(voir aussi Morphismes, Lemmes \ref{morphisms-lemma-open-immersion-etale}
et \ref{morphisms-lemma-composition-etale}).
\end{proof}














\section{Présentations d'espaces algébriques}
\label{section-presentations}

\noindent
Étant donné un espace algébrique, on peut en trouver une «~présentation~».


\begin{lemma}
\label{lemma-space-presentation}
Soit $F$ un espace algébrique au-dessus de $S$. Soit $f : U \to F$ un
morphisme \'etale surjectif d'un schéma vers $F$. Posons $R = U \times_F U$.
Alors
\begin{enumerate}
\item $j : R \to U \times_S U$ définit sur
$U$ une relation d'équivalence au-dessus de $S$ (voir
Groupoïdes, Définition \ref{groupoids-definition-equivalence-relation}).
\item les morphismes $s, t : R \to U$ sont \'etales, et
\item le diagramme
$$
\xymatrix{
R \ar@<1ex>[r] \ar@<-1ex>[r] &
U \ar[r] &
F
}
$$
est un diagramme de conoyau dans $\Sh((\Sch/S)_{fppf})$.
\end{enumerate}
\end{lemma}

\begin{proof}
Soit $T/S$ un objet de $(\Sch/S)_{fppf}$.
Alors $R(T) = \{(a, b) \in U(T) \times U(T) \mid f \circ a = f \circ b\}$,
ce qui définit une relation d'équivalence sur $U(T)$.
Les morphismes $s, t : R \to U$ sont \'etales parce que le morphisme
$U \to F$ est \'etale.

\medskip\noindent
Pour démontrer (3), montrons d'abord que
$U \to F$ est un morphisme surjectif de faisceaux; voir
Sites, Définition \ref{sites-definition-sheaves-injective-surjective}.
Soit $\xi \in F(T)$, où $T$ est comme ci-dessus. Posons $V = T \times_{\xi, F, f}U$.
Par hypothèse, $V$ est un schéma et $V \to T$ est un morphisme \'etale surjectif.
Ainsi, $\{V \to T\}$ est un recouvrement pour la topologie fppf.
Puisque, par construction, $\xi|_V$ se factorise par $U$, on
en déduit que $U \to F$ est surjectif. La surjectivité implique que
$F$ est le conoyau du diagramme d'après
Sites, Lemme \ref{sites-lemma-coequalizer-surjection}.
\end{proof}

\noindent
Ce lemme conduit aux définitions suivantes.

\begin{definition}
\label{definition-etale-equivalence-relation}
Soit $S$ un schéma. Soit $U$ un schéma au-dessus de $S$.
Une {\it relation d'équivalence \'etale} sur $U$ au-dessus de $S$
est une relation d'équivalence $j : R \to U \times_S U$
telle que $s, t : R \to U$ soient des morphismes \'etales de schémas.
\end{definition}

\begin{definition}
\label{definition-presentation}
Soit $F$ un espace algébrique au-dessus de $S$.
Une {\it présentation} de $F$ est la donnée d'un schéma
$U$ au-dessus de $S$, d'une relation d'équivalence \'etale $R$ sur $U$ au-dessus de $S$, et
d'un morphisme \'etale surjectif $U \to F$ tel que $R = U \times_F U$.
\end{definition}

\noindent
De manière équivalente, on pourrait demander l'existence d'un isomorphisme
$$
U/R \cong F
$$
où le quotient $U/R$ est défini comme dans
Groupoïdes, Section \ref{groupoids-section-quotient-sheaves}.
Pour construire des espaces algébriques, nous étudierons la question réciproque, à savoir
pour quelles relations d'équivalence le faisceau quotient $U/R$ est un espace algébrique.
On verra finalement qu'il en est toujours ainsi si $R$ est une relation
d'équivalence \'etale sur $U$ au-dessus de $S$; voir le Théorème \ref{theorem-presentation}.





























\section{Espaces algébriques et relations d'équivalence}
\label{section-spaces-from-equivalence-relations}

\noindent
Supposons donnés un schéma $U$ au-dessus de $S$
et une relation d'équivalence \'etale $R$ sur $U$ au-dessus de $S$.
Nous voudrions montrer que ces données définissent un espace algébrique.
Nous allons établir une suite de lemmes montrant que le faisceau quotient $U/R$
(voir Groupoïdes, Définition \ref{groupoids-definition-quotient-sheaf})
possède toutes les propriétés que lui impose la
Définition \ref{definition-algebraic-space}.

\begin{lemma}
\label{lemma-pullback-etale-equivalence-relation}
Soit $S$ un schéma. Soit $U$ un schéma au-dessus de $S$.
Soit $j = (s, t) : R \to U \times_S U$
une relation d'équivalence \'etale sur $U$ au-dessus de $S$.
Soit $U' \to U$ un morphisme \'etale.
Soit $R'$ la restriction de $R$ à $U'$; voir
Groupoïdes, Définition \ref{groupoids-definition-restrict-relation}.
Alors $j' : R' \to U' \times_S U'$ est également une relation
d'équivalence \'etale.
\end{lemma}

\begin{proof}
Il résulte de la description de $s', t'$ dans
Groupoïdes, Lemme \ref{groupoids-lemma-restrict-groupoid}
que $s' , t' : R' \to U'$ sont \'etales,
car ce sont des composés de changements de base de morphismes \'etales
(voir Morphismes, Lemme \ref{morphisms-lemma-base-change-etale}
et \ref{morphisms-lemma-composition-etale}).
\end{proof}

\noindent
Nous utiliserons souvent le lemme suivant pour trouver des sous-espaces ouverts des espaces
algébriques. Une légère amélioration de ce lemme (sous des hypothèses plus générales) est
Amorçage, Lemme \ref{bootstrap-lemma-better-finding-opens}.

\begin{lemma}
\label{lemma-finding-opens}
Soit $S$ un schéma.
Soit $U$ un schéma au-dessus de $S$.
Soit $j = (s, t) : R \to U \times_S U$ une pré-relation.
Soit $g : U' \to U$ un morphisme.
Supposons que
\begin{enumerate}
\item $j$ est une relation d'équivalence,
\item $s, t : R \to U$ sont surjectifs, plats et
localement de présentation finie,
\item $g$ est plat et localement de présentation finie.
\end{enumerate}
Soit $R' = R|_{U'}$ la restriction de $R$ à $U'$. Alors
$U'/R' \to U/R$ est représentable et est une immersion ouverte.
\end{lemma}

\begin{proof}
D'après Groupoïdes, Lemme \ref{groupoids-lemma-restrict-relation},
le morphisme $j' = (s', t') : R' \to U' \times_S U'$
définit une relation d'équivalence. Puisque $g$ est plat et localement de
présentation finie, $g$ est aussi universellement ouvert
(Morphismes, Lemme \ref{morphisms-lemma-fppf-open}).
Pour la même raison, $s, t$ sont également universellement ouverts.
Posons $W^1 = g(U') \subset U$ et $W = t(s^{-1}(W^1))$.
Alors $W^1$ et $W$ sont ouverts dans $U$. De plus, puisque $j$ est une
relation d'équivalence, on a $t(s^{-1}(W)) = W$ (voir
Groupoïdes, Lemme \ref{groupoids-lemma-constructing-invariant-opens},
par exemple).

\medskip\noindent
D'après
Groupoïdes,
Lemme \ref{groupoids-lemma-quotient-pre-equivalence-relation-restrict},
le morphisme de faisceaux $F' = U'/R' \to F = U/R$ est injectif.
Soit $a : T \to F$ un morphisme d'un schéma vers $U/R$.
Il faut montrer que $T \times_F F'$ est représentable
par un sous-schéma ouvert de $T$.

\medskip\noindent
Le morphisme $a$ est donné par les données suivantes :
un recouvrement fppf $\{\varphi_j : T_j \to T\}_{j \in J}$ de $T$ et des
morphismes $a_j : T_j \to U$ tels que les morphismes
$$
a_j \times a_{j'} :
T_j \times_T T_{j'}
\longrightarrow
U \times_S U
$$
se factorisent par $j : R \to U \times_S U$ au moyen de morphismes (uniques)
$r_{jj'} : T_j \times_T T_{j'} \to R$. Le système
$(a_j)$ correspond à $a$ en ce sens que les diagrammes
$$
\xymatrix{
T_j \ar[r]_{a_j} \ar[d] & U \ar[d] \\
T \ar[r]^a & F
}
$$
commutent.

\medskip\noindent
Considérons les ouverts $W_j = a_j^{-1}(W) \subset T_j$.
Puisque $t(s^{-1}(W)) = W$, on voit que
$$
W_j \times_T T_{j'} =
r_{jj'}^{-1}(t^{-1}(W)) = r_{jj'}^{-1}(s^{-1}(W)) =
T_j \times_T W_{j'}.
$$
D'après
Descente, Lemme \ref{descent-lemma-open-fpqc-covering},
cela signifie qu'il existe un ouvert
$W_T \subset T$ tel que $\varphi_j^{-1}(W_T) = W_j$ pour tout $j \in J$.
Nous affirmons que $W_T \to T$ représente $T \times_F F' \to T$.

\medskip\noindent
Montrons d'abord que $W_T \to T \to F$ définit un élément de
$F'(W_T)$. Puisque $\{W_j \to W_T\}_{j \in J}$ est un
recouvrement fppf de $W_T$, il suffit de montrer que
chaque $W_j \to U \to F$ définit un élément de $F'(W_j)$ (car $F'$ est un faisceau
pour la topologie fppf). Considérons le diagramme commutatif
$$
\xymatrix{
W'_j \ar[rr] \ar[dd] \ar[rd] & & U' \ar[d]^g \\
& s^{-1}(W^1) \ar[r]_s \ar[d]^t & W^1 \ar[d] \\
W_j \ar[r]^{a_j|_{W_j}} & W \ar[r] & F
}
$$
où $W'_j = W_j \times_W s^{-1}(W^1) \times_{W^1} U'$.
Puisque $t$ et $g$ sont surjectifs, plats et localement de présentation
finie, il en est de même de $W'_j \to W_j$. Par conséquent, la restriction de
l'élément $W_j \to U \to F$ à $W'_j$ est un élément de $F'$,
comme voulu.

\medskip\noindent
Supposons que $f : T' \to T$ soit un morphisme de schémas
tel que $a|_{T'} \in F'(T')$. Il faut montrer que
$f$ se factorise par l'ouvert $W_T$. Puisque
$\{T' \times_T T_j \to T'\}$ est un recouvrement fppf de $T'$,
il suffit de montrer que chaque $T' \times_T T_j \to T$
se factorise par $W_T$. Nous pouvons donc supposer que $f$ se factorise
sous la forme $\varphi_j \circ f_j : T' \to T_j \to T$ pour un certain $j$.
Dans ce cas, la condition $a|_{T'} \in F'(T')$ signifie qu'il existe
un recouvrement fppf $\{\psi_i : T'_i \to T'\}_{i \in I}$ et des
morphismes $b_i : T'_i \to U'$ tels que
$$
\xymatrix{
T'_i \ar[r]_{b_i} \ar[d]_{f_j \circ \psi_i} & U' \ar[r]_g & U \ar[d] \\
T_j \ar[r]^{a_j} & U \ar[r] & F
}
$$
soit commutatif. Cette commutativité signifie qu'il existe un
morphisme $r'_i : T'_i \to R$ tel que
$t \circ r'_i = a_j \circ f_j \circ \psi_i$ et
$s \circ r'_i = g \circ b_i$. Il en résulte que
$\Im(f_j \circ \psi_i) \subset W_j$, ce qui conclut.
\end{proof}

\noindent
Le lemme suivant n'est pas tout à fait trivial, bien qu'il semble
devoir l'être.

\begin{lemma}
\label{lemma-when-it-works-it-works}
Soit $S$ un schéma. Soit $U$ un schéma au-dessus de $S$.
Soit $j = (s, t) : R \to U \times_S U$
une relation d'équivalence \'etale sur $U$ au-dessus de $S$.
Si le quotient $U/R$ est un espace algébrique, alors
$U \to U/R$ est \'etale et surjectif. Par conséquent,
$(U, R, U \to U/R)$ est une présentation de l'espace
algébrique $U/R$.
\end{lemma}

\begin{proof}
Notons $c : U \to U/R$ le morphisme en question.
Soit $T$ un schéma et soit $a : T \to U/R$ un morphisme.
Il faut montrer que le morphisme (de schémas)
$\pi : T \times_{a, U/R, c} U \to T$ est \'etale et surjectif.
Le morphisme $a$ correspond à un recouvrement fppf
$\{\varphi_i : T_i \to T\}$ et à des morphismes $a_i : T_i \to U$ tels
que
$a_i \times a_{i'} : T_i \times_T T_{i'} \to U \times_S U$
se factorise par $R$, et tels que $c \circ a_i = a \circ \varphi_i$.
Ainsi,
$$
T_i \times_{\varphi_i, T} T \times_{a, U/R, c} U =
T_i \times_{c \circ a_i, U/R, c} U =
T_i \times_{a_i, U} U \times_{c, U/R, c} U = T_i \times_{a_i, U, t} R.
$$
Puisque $t$ est \'etale et surjectif, on en déduit que
le changement de base de $\pi$ à $T_i$ est surjectif et \'etale.
Comme la propriété d'être surjectif et \'etale est locale
sur la base pour la topologie fpqc (voir la
Remarque \ref{remark-list-properties-fpqc-local-base}),
cela conclut.
\end{proof}

\begin{lemma}
\label{lemma-presentation-quasi-compact}
Soit $S$ un schéma.
Soit $U$ un schéma au-dessus de $S$.
Soit $j = (s, t) : R \to U \times_S U$
une relation d'équivalence \'etale sur $U$ au-dessus de $S$.
Supposons $U$ affine. Alors le quotient $F = U/R$
est un espace algébrique, et $U \to F$ est \'etale et surjectif.
\end{lemma}

\begin{proof}
Puisque $j : R \to U \times_S U$ est un monomorphisme, $j$ est séparé
(voir Schémas, Lemme \ref{schemes-lemma-monomorphism-separated}).
Comme $U$ est affine, on voit que $U \times_S U$
(qui est muni d'un monomorphisme vers le schéma affine
$U \times U$) est séparé. Il s'ensuit que $R$ est séparé.
En particulier, les morphismes $s, t$ sont séparés et \'etales.

\medskip\noindent
Puisque le composé $R \to U \times_S U \to U$ est
localement de type fini, on en déduit que
$j$ est localement de type fini (voir
Morphismes, Lemme \ref{morphisms-lemma-permanence-finite-type}).
Comme $j$ est aussi un monomorphisme, ses fibres sont finies, et
l'on voit que $j$ est localement quasi-fini d'après
Morphismes, Lemme \ref{morphisms-lemma-finite-fibre}.
En définitive, $j$ est séparé et localement quasi-fini.

\medskip\noindent
La première étape consiste à montrer que le morphisme de passage au quotient
$c : U \to F$ est représentable.
Considérons un schéma $T$ et un morphisme $a : T \to F$.
Il faut montrer que le faisceau $G = T \times_{a, F, c} U$
est représentable.
Comme on l'a vu dans les démonstrations des Lemmes \ref{lemma-finding-opens} et
\ref{lemma-when-it-works-it-works}, il existe un recouvrement fppf
$\{\varphi_i : T_i \to T\}_{i \in I}$ et des morphismes $a_i : T_i \to U$
tels que $a_i \times a_{i'} : T_i \times_T T_{i'} \to U \times_S U$
se factorise par $R$, et tels que $c \circ a_i = a \circ \varphi_i$.
Comme dans la démonstration du Lemme \ref{lemma-when-it-works-it-works}, on voit que
\begin{eqnarray*}
T_i \times_{\varphi_i, T} G & = &
T_i \times_{\varphi_i, T} T \times_{a, U/R, c} U \\
& = & T_i \times_{c \circ a_i, U/R, c} U \\
& = & T_i \times_{a_i, U} U \times_{c, U/R, c} U \\
& = & T_i \times_{a_i, U, t} R
\end{eqnarray*}
Puisque $t$ est séparé et \'etale, donc en particulier
séparé et localement quasi-fini (d'après Morphismes, Lemmes
\ref{morphisms-lemma-unramified-quasi-finite} et
\ref{morphisms-lemma-flat-unramified-etale}),
on voit que la restriction
de $G$ à chaque $T_i$ est représentable par un morphisme de schémas
$X_i \to T_i$ séparé et localement quasi-fini. D'après
Descente, Lemme \ref{descent-lemma-descent-data-sheaves},
on obtient une donnée de descente $(X_i, \varphi_{ii'})$ relativement
au recouvrement fppf $\{T_i \to T\}$. Comme chaque
$X_i \to T_i$ est séparé et localement quasi-fini, il résulte de
Compléments sur les morphismes, Lemme
\ref{more-morphisms-lemma-separated-locally-quasi-finite-morphisms-fppf-descend}
que cette donnée de descente est effective.
Ainsi, d'après
Descente, Lemme \ref{descent-lemma-descent-data-sheaves} (2),
on en déduit que $G$ est représentable, comme voulu.

\medskip\noindent
La deuxième étape de la démonstration consiste à montrer que $U \to F$ est surjectif et
\'etale. Cela résulte de ce qui précède : dans la première étape, nous avons
vu que $G = T \times_{a, F, c} U$ est un schéma au-dessus de $T$ dont les changements de base
sont les schémas $X_i \to T_i$, qui sont surjectifs et \'etales. Ainsi, $G \to T$
est surjectif et \'etale (voir la
Remarque \ref{remark-list-properties-fpqc-local-base}).
On peut également reprendre la démonstration du
Lemme \ref{lemma-when-it-works-it-works} dans la présente
situation.

\medskip\noindent
La troisième et dernière étape consiste à montrer que le morphisme diagonal $F \to F \times F$
est représentable. Observons d'abord que le diagramme
$$
\xymatrix{
R \ar[r] \ar[d]_j & F \ar[d]^\Delta \\
U \times_S U \ar[r] & F \times F
}
$$
est un carré cartésien. D'après le
Lemme \ref{lemma-product-representable-transformations}, le morphisme
$U \times_S U \to F \times F$ est représentable (notons que
$h_U \times h_U = h_{U \times_S U}$). De plus, d'après le
Lemme \ref{lemma-product-representable-transformations-property},
le morphisme $U \times_S U \to F \times F$ est surjectif
et \'etale (notons aussi que \'etale et surjectif figurent dans les listes des
Remarques \ref{remark-list-properties-fpqc-local-base}
et \ref{remark-list-properties-stable-composition}).
Il résulte soit du
Lemme \ref{lemma-base-change-representable-transformations}
et du diagramme ci-dessus, soit de l'écriture de $R \to F$ comme $R \to U \to F$ et des
Lemmes
\ref{lemma-morphism-schemes-gives-representable-transformation} et
\ref{lemma-composition-representable-transformations}, que
$R \to F$ est également représentable. Soient $T$ un schéma et
$a : T \to F \times F$ un morphisme. Il faut montrer que
$G = T \times_{a, F \times F, \Delta} F$ est représentable.
D'après ce qui précède, le morphisme (de schémas)
$$
T' = (U \times_S U) \times_{F \times F, a} T \longrightarrow T
$$
est surjectif et \'etale. Ainsi, $\{T' \to T\}$ est un recouvrement
\'etale de $T$. Notons également que
$$
T' \times_T G = T' \times_{U \times_S U, j} R
$$
comme on le voit en examinant le cube suivant
$$
\xymatrix{
& R \ar[rr] \ar[dd] & & F \ar[dd] \\
T' \times_T G \ar[rr] \ar[dd] \ar[ru] & & G \ar[dd] \ar[ru] & \\
& U \times_S U \ar'[r][rr] & & F \times F \\
T' \ar[rr] \ar[ru] & & T \ar[ru]
}
$$
On voit ainsi que la restriction de $G$ à $T'$ est représentable
par un schéma $X$, et de plus que le morphisme $X \to T'$ est
un changement de base du morphisme $j$. Par conséquent, $X \to T'$ est
séparé et localement quasi-fini (voir le deuxième paragraphe de la démonstration).
D'après Descente, Lemme \ref{descent-lemma-descent-data-sheaves},
on obtient une donnée de descente $(X, \varphi)$ relativement
au recouvrement fppf $\{T' \to T\}$. Puisque
$X \to T'$ est séparé et localement quasi-fini, il résulte de
Compléments sur les morphismes, Lemme
\ref{more-morphisms-lemma-separated-locally-quasi-finite-morphisms-fppf-descend}
que cette donnée de descente est effective.
Ainsi, d'après
Descente, Lemme \ref{descent-lemma-descent-data-sheaves} (2),
on en déduit que $G$ est représentable, comme voulu.
\end{proof}

\begin{theorem}
\label{theorem-presentation}
Soit $S$ un schéma. Soit $U$ un schéma au-dessus de $S$.
Soit $j = (s, t) : R \to U \times_S U$
une relation d'équivalence \'etale sur $U$ au-dessus de $S$.
Alors le quotient $U/R$ est un espace algébrique,
et $U \to U/R$ est \'etale et surjectif; autrement dit,
$(U, R, U \to U/R)$ est une présentation de $U/R$.
\end{theorem}

\begin{proof}
D'après le Lemme \ref{lemma-when-it-works-it-works},
il suffit de prouver que $U/R$ est un espace algébrique.
Soit $U' \to U$ un morphisme surjectif et \'etale.
Alors $\{U' \to U\}$ est en particulier un recouvrement fppf.
Soit $R'$ la restriction de $R$ à $U'$; voir
Groupoïdes, Définition \ref{groupoids-definition-restrict-relation}.
D'après
Groupoïdes, Lemme \ref{groupoids-lemma-quotient-groupoid-restrict},
on a $U/R \cong U'/R'$.
D'après le Lemme \ref{lemma-pullback-etale-equivalence-relation}, $R'$ est une
relation d'équivalence \'etale sur $U'$. Nous pouvons donc remplacer $U$ par $U'$.

\medskip\noindent
Appliquons la remarque précédente à $U' = \coprod U_i$, où
$U = \bigcup U_i$ est un recouvrement ouvert affine de $U$. Nous
pouvons donc supposer, et nous le ferons, que $U = \coprod U_i$, où
chaque $U_i$ est un schéma affine.

\medskip\noindent
Considérons la restriction $R_i$ de $R$ à $U_i$.
D'après le Lemme \ref{lemma-pullback-etale-equivalence-relation},
c'est une relation d'équivalence \'etale.
Posons $F_i = U_i/R_i$ et $F = U/R$.
Il est clair que $\coprod F_i \to F$ est surjectif.
D'après le Lemme \ref{lemma-finding-opens}, chaque $F_i \to F$
est représentable et est une immersion ouverte.
En appliquant le Lemme \ref{lemma-presentation-quasi-compact}
à $(U_i, R_i)$, on voit que $F_i$ est un espace algébrique.
Puis, d'après le Lemme \ref{lemma-when-it-works-it-works},
$U_i \to F_i$ est \'etale et surjectif.
Il résulte du Lemme \ref{lemma-coproduct-algebraic-spaces}
que $\coprod F_i$ est un espace algébrique.
Enfin, nous avons vérifié toutes les
hypothèses du Lemme \ref{lemma-glueing-algebraic-spaces},
et il s'ensuit que $F = U/R$ est un espace algébrique.
\end{proof}










\section{Espaces algébriques, reformulation}
\label{section-algebraic-spaces-retrofitted}

\noindent
Nous commençons à constituer notre arsenal de lemmes sur les espaces algébriques.
Le premier résultat affirme que, dans la Définition \ref{definition-algebraic-space},
on peut affaiblir comme suit la condition portant sur la diagonale.

\begin{lemma}
\label{lemma-etale-locally-representable-gives-space}
Soit $S$ un schéma appartenant à $\Sch_{fppf}$.
Soit $F$ un faisceau sur $(\Sch/S)_{fppf}$
tel qu'il existe $U \in \Ob((\Sch/S)_{fppf})$ et un morphisme
$U \to F$ représentable, surjectif et \'etale.
Alors $F$ est un espace algébrique.
\end{lemma}

\begin{proof}
Posons $R = U \times_F U$. C'est un schéma, car $U \to F$ est supposé représentable.
Les projections $s, t : R \to U$ sont \'etales, car $U \to F$ est supposé \'etale.
Le morphisme $j = (t, s) : R \to U \times_S U$ est un monomorphisme et définit une relation
d'équivalence, puisque $R = U \times_F U$. D'après le Théorème \ref{theorem-presentation},
le faisceau quotient $F' = U/R$ est un espace algébrique et $U \to F'$
est surjectif et \'etale. À nouveau, comme $R = U \times_F U$, on obtient
une factorisation canonique $U \to F' \to F$, et $F' \to F$ est un morphisme injectif
de faisceaux. D'autre part, $U \to F$ est surjectif comme morphisme
de faisceaux d'après le Lemme \ref{lemma-surjective-flat-locally-finite-presentation}.
Ainsi, $F' \to F$ est également surjectif, et l'on conclut que $F' = F$ est un
espace algébrique.
\end{proof}

\begin{lemma}
\label{lemma-etale-locally-representable-by-space-gives-space}
Soit $S$ un schéma appartenant à $\Sch_{fppf}$. Soit $G$ un espace
algébrique au-dessus de $S$, soit $F$ un faisceau sur $(\Sch/S)_{fppf}$, et soit
$G \to F$ un morphisme représentable de foncteurs qui est
surjectif et \'etale. Alors $F$ est un espace algébrique.
\end{lemma}

\begin{proof}
Choisissons un schéma $U$ et un morphisme surjectif \'etale $U \to G$.
Puisque $G$ est un espace algébrique, $U \to G$ est représentable.
Le composé $U \to G \to F$ est donc représentable,
surjectif et \'etale. Voir les Lemmes
\ref{lemma-composition-representable-transformations} and
\ref{lemma-composition-representable-transformations-property}.
Ainsi, $F$ est un espace algébrique d'après le
Lemme \ref{lemma-etale-locally-representable-gives-space}.
\end{proof}

\begin{lemma}
\label{lemma-representable-over-space}
\begin{slogan}
Un foncteur qui admet un morphisme représentable vers un espace algébrique est
un espace algébrique.
\end{slogan}
Soit $S$ un schéma appartenant à $\Sch_{fppf}$.
Soit $F$ un espace algébrique au-dessus de $S$.
Soit $G \to F$ un morphisme représentable de foncteurs.
Alors $G$ est un espace algébrique.
\end{lemma}

\begin{proof}
D'après le Lemme \ref{lemma-representable-transformation-to-sheaf},
on voit que $G$ est un faisceau. Le diagramme
$$
\xymatrix{
G \times_F G \ar[r] \ar[d] & F \ar[d]^{\Delta_F} \\
G \times G \ar[r] & F \times F
}
$$
est cartésien. On voit donc que $G \times_F G \to G \times G$
est représentable d'après le
Lemme \ref{lemma-base-change-representable-transformations}.
D'après le
Lemme \ref{lemma-representable-transformation-diagonal},
on voit que $G \to G \times_F G$ est représentable.
Ainsi, $\Delta_G : G \to G \times G$ est représentable comme composé
de morphismes représentables de foncteurs; voir le
Lemme \ref{lemma-composition-representable-transformations}.
Enfin, soit $U$ un objet de $(\Sch/S)_{fppf}$
et soit $U \to F$ surjectif et \'etale. Par hypothèse,
$U \times_F G$ est représentable par un schéma $U'$. D'après le
Lemme \ref{lemma-base-change-representable-transformations-property},
le morphisme $U' \to G$ est surjectif et \'etale. Cela vérifie
la dernière condition de la Définition \ref{definition-algebraic-space}, ce qui conclut.
\end{proof}

\begin{lemma}
\label{lemma-representable-morphisms-spaces-property}
Soit $S$ un schéma appartenant à $\Sch_{fppf}$.
Soient $F$, $G$ des espaces algébriques au-dessus de $S$.
Soit $G \to F$ un morphisme représentable.
Soient $U \in \Ob((\Sch/S)_{fppf})$ et $q : U \to F$
surjectif et \'etale. Posons $V = G \times_F U$.
Enfin, soit $\mathcal{P}$ une propriété des morphismes
de schémas comme dans la Définition \ref{definition-relative-representable-property}.
Alors $G \to F$ possède la propriété $\mathcal{P}$ si et seulement si
$V \to U$ possède la propriété $\mathcal{P}$.
\end{lemma}

\begin{proof}
(Ce lemme résulte des
Lemmes \ref{lemma-base-change-representable-transformations-property} et
\ref{lemma-descent-representable-transformations-property},
mais nous en donnons aussi ici une démonstration directe.)
Il résulte immédiatement des définitions que, si $G \to F$ possède la propriété
$\mathcal{P}$, alors $V \to U$ possède la propriété $\mathcal{P}$.
Réciproquement, supposons que $V \to U$ possède la propriété $\mathcal{P}$.
Soit $T \to F$ un morphisme d'un schéma vers $F$.
Soit $T' = T \times_F G$, qui est un schéma puisque $G \to F$ est
représentable. Il faut montrer que $T' \to T$ possède la propriété $\mathcal{P}$.
Considérons le diagramme commutatif de schémas
$$
\xymatrix{
V \ar[d] & T \times_F V \ar[d] \ar[l] \ar[r] &
T \times_F G \ar[d] \ar@{=}[r] & T' \\
U & T \times_F U \ar[l] \ar[r] & T
}
$$
où les deux carrés sont cartésiens. On en déduit que
la flèche médiane possède la propriété $\mathcal{P}$ en tant que changement de base
de $V \to U$. Enfin, $\{T \times_F U \to T\}$ est un recouvrement fppf,
car il est surjectif \'etale; on en déduit donc que
$T' \to T$ possède la propriété $\mathcal{P}$, puisque celle-ci est locale sur la
base pour la topologie fppf.
\end{proof}

\begin{lemma}
\label{lemma-morphism-sheaves-with-P-effective-descent-etale}
Let $S$ be a scheme contained in $\Sch_{fppf}$.
Let $G \to F$ be a transformation of presheaves on $(\Sch/S)_{fppf}$.
Let $\mathcal{P}$ be a property of morphisms of schemes.
Assume
\begin{enumerate}
\item $\mathcal{P}$ is preserved under any base change, fppf local on the
base, and morphisms of type $\mathcal{P}$ satisfy descent for fppf coverings,
see Descent, Definition \ref{descent-definition-descending-types-morphisms},
\item $G$ is a sheaf,
\item $F$ is an algebraic space,
\item there exists a $U \in \Ob((\Sch/S)_{fppf})$
and a surjective \'etale morphism $U \to F$ such that
$V = G \times_F U$ is representable, and
\item $V \to U$ has $\mathcal{P}$.
\end{enumerate}
Then $G$ is an algebraic space, $G \to F$ is representable and has property
$\mathcal{P}$.
\end{lemma}

\begin{proof}
Let $T$ be a scheme and let $T \to F$ be a morphism. Then
$U \times_F T \to T$ is surjective \'etale, hence $\{U \times_F T \to T\}$ is
a covering for the \'etale topology. Consider
$$
W = G \times_F (U \times_F T) = V \times_F T = V \times_U (U \times_F T).
$$
It is a scheme since $F$ is an algebraic space. The morphism
$W \to U \times_F T$ has property $\mathcal{P}$ since it is a
base change of $V \to U$. There is an isomorphism
\begin{align*}
W \times_T (U \times_F T) & =
(G \times_F (U \times_F T)) \times_T (U \times_F T) \\
& = (U \times_F T) \times_T (G \times_F (U \times_F T)) \\
& = (U \times_F T) \times_T W
\end{align*}
over $(U \times_F T) \times_T (U \times_F T)$. The middle equality maps
$((g, (u_1, t)), (u_2, t))$ to $((u_1, t), (g, (u_2, t)))$.
This defines a descent datum for $W/U \times_F T/T$, see
Descent, Definition \ref{descent-definition-descent-datum}.
This follows from
Descent, Lemma \ref{descent-lemma-descent-data-sheaves}.
Namely we have a sheaf $G \times_F T$, whose
base change to $U \times_F T$ is represented by $W$ and the isomorphism
above is the one from the proof of
Descent, Lemma \ref{descent-lemma-descent-data-sheaves}.
By assumption on $\mathcal{P}$ the descent datum above is representable.
Hence by the last statement of
Descent, Lemma \ref{descent-lemma-descent-data-sheaves}
we see that $G \times_F T$ is representable. This proves that
$G \to F$ is a representable transformation of functors.

\medskip\noindent
As $G \to F$ is representable, we see that $G$ is an algebraic space by
Lemma \ref{lemma-representable-over-space}. The fact that $G \to F$ has
property $\mathcal{P}$ now follows from
Lemma \ref{lemma-representable-morphisms-spaces-property}.
\end{proof}

\begin{lemma}
\label{lemma-lift-morphism-presentations}
Let $S$ be a scheme contained in $\Sch_{fppf}$.
Let $F, G$ be algebraic spaces over $S$.
Let $a : F \to G$ be a morphism.
Given any $V \in \Ob((\Sch/S)_{fppf})$
and a surjective \'etale morphism $q : V \to G$ there exists
a $U \in \Ob((\Sch/S)_{fppf})$
and a commutative diagram
$$
\xymatrix{
U \ar[d]_p \ar[r]_\alpha &
V \ar[d]^q \\
F \ar[r]^a & G
}
$$
with $p$ surjective and \'etale.
\end{lemma}

\begin{proof}
First choose $W \in \Ob((\Sch/S)_{fppf})$
with surjective \'etale morphism $W \to F$.
Next, put $U = W \times_G V$. Since $G$ is an algebraic space
we see that $U$ is isomorphic to an object of $(\Sch/S)_{fppf}$.
As $q$ is surjective \'etale, we see that $U \to W$ is surjective
\'etale (see
Lemma \ref{lemma-base-change-representable-transformations-property}).
Thus $U \to F$ is surjective \'etale as a composition of surjective
\'etale morphisms (see
Lemma \ref{lemma-composition-representable-transformations-property}).
\end{proof}



























\section{Immersions and Zariski coverings of algebraic spaces}
\label{section-Zariski}

\noindent
At this point an interesting phenomenon occurs. We have already defined
the notion of an open immersion of algebraic spaces (through
Definition \ref{definition-relative-representable-property})
but we have yet to define the notion of a {\it point}\footnote{We
will associate a topological space to an algebraic space in
Properties of Spaces, Section
\ref{spaces-properties-section-points},
and its opens will correspond exactly to the open subspaces defined below.}.
Thus the {\it Zariski topology} of an algebraic space
has already been defined, but there is no space yet!

\medskip\noindent
Perhaps superfluously we formally introduce immersions as follows.

\begin{definition}
\label{definition-immersion}
Let $S \in \Ob(\Sch_{fppf})$ be a scheme.
Let $F$ be an algebraic space over $S$.
\begin{enumerate}
\item A morphism of algebraic spaces over $S$
is called an {\it open immersion} if it is representable, and an open immersion
in the sense of Definition \ref{definition-relative-representable-property}.
\item An {\it open subspace} of $F$ is a subfunctor $F' \subset F$
such that $F'$ is an algebraic space and $F' \to F$ is an
open immersion.
\item A morphism of algebraic spaces over $S$
is called a {\it closed immersion} if it is representable, and a closed
immersion in the sense of
Definition \ref{definition-relative-representable-property}.
\item A {\it closed subspace} of $F$ is a subfunctor $F' \subset F$
such that $F'$ is an algebraic space and $F' \to F$ is a
closed immersion.
\item A morphism of algebraic spaces over $S$
is called an {\it immersion} if it is representable, and an immersion
in the sense of Definition \ref{definition-relative-representable-property}.
\item A {\it locally closed subspace} of $F$ is a subfunctor $F' \subset F$
such that $F'$ is an algebraic space and $F' \to F$ is an
immersion.
\end{enumerate}
\end{definition}

\noindent
We note that these definitions make sense since an immersion
is in particular a monomorphism (see
Schemes, Lemma \ref{schemes-lemma-immersions-monomorphisms}
and Lemma \ref{lemma-representable-transformations-property-implication}),
and hence the image of an
immersion $G \to F$ of algebraic spaces is a subfunctor $F' \subset F$
which is (canonically) isomorphic to $G$. Thus some of the discussion
of Schemes, Section \ref{schemes-section-immersions} carries over to the
setting of algebraic spaces.

\begin{lemma}
\label{lemma-composition-immersions}
Let $S \in \Ob(\Sch_{fppf})$ be a scheme.
A composition of (closed, resp.\ open) immersions of
algebraic spaces over $S$ is a (closed, resp.\ open)
immersion of algebraic spaces over $S$.
\end{lemma}

\begin{proof}
See Lemma \ref{lemma-composition-representable-transformations-property} and
Remarks \ref{remark-list-properties-fpqc-local-base} (see very last line of
that remark) and \ref{remark-list-properties-stable-composition}.
\end{proof}

\begin{lemma}
\label{lemma-base-change-immersions}
Let $S \in \Ob(\Sch_{fppf})$ be a scheme.
A base change of a (closed, resp.\ open) immersion
of algebraic spaces over $S$ is a (closed, resp.\ open)
immersion of algebraic spaces over $S$.
\end{lemma}

\begin{proof}
See Lemma \ref{lemma-base-change-representable-transformations-property} and
Remark \ref{remark-list-properties-fpqc-local-base} (see very last line of
that remark).
\end{proof}

\begin{lemma}
\label{lemma-sub-subspaces}
Let $S \in \Ob(\Sch_{fppf})$ be a scheme.
Let $F$ be an algebraic space over $S$. Let $F_1$, $F_2$ be
locally closed subspaces of $F$. If $F_1 \subset F_2$ as subfunctors
of $F$, then $F_1$ is a locally closed subspace of $F_2$.
Similarly for closed and open subspaces.
\end{lemma}

\begin{proof}
Let $T \to F_2$ be a morphism with $T$ a scheme.
Since $F_2 \to F$ is a monomorphism, we see that
$T \times_{F_2} F_1 = T \times_F F_1$. The lemma follows
formally from this.
\end{proof}

\noindent
Let us formally define the notion of a Zariski open covering of
algebraic spaces. Note that in Lemma \ref{lemma-glueing-algebraic-spaces}
we have already encountered such open coverings as a method for
constructing algebraic spaces.

\begin{definition}
\label{definition-Zariski-open-covering}
Let $S \in \Ob(\Sch_{fppf})$ be a scheme.
Let $F$ be an algebraic space over $S$.
A {\it Zariski covering} $\{F_i \subset F\}_{i \in I}$ of $F$
is given by a set $I$ and a collection of open subspaces
$F_i \subset F$ such that $\coprod F_i \to F$ is a surjective
map of sheaves.
\end{definition}

\noindent
Note that if $T$ is a schemes,
and $a : T \to F$ is a morphism, then each of the fibre products
$T \times_F F_i$ is identified with an open subscheme
$T_i \subset T$. The final condition of the definition signifies
exactly that $T = \bigcup_{i \in I} T_i$.

\medskip\noindent
It is clear that the collection $F_{Zar}$ of open subspaces of
$F$ is a set (as $(\Sch/S)_{fppf}$ is a site, hence a set).
Moreover, we can turn $F_{Zar}$ into a category by letting the
morphisms be inclusions of subfunctors (which are automatically open
immersions by Lemma \ref{lemma-sub-subspaces}). Finally,
Definition \ref{definition-Zariski-open-covering}
provides the notion of a Zariski covering $\{F_i \to F'\}_{i \in I}$
in the category $F_{Zar}$. Hence, just as in the case of a topological
space (see Sites, Example \ref{sites-example-site-topological})
by suitably choosing a set of coverings
we may obtain a Zariski site of the algebraic space $F$.

\begin{definition}
\label{definition-small-Zariski-site}
Let $S \in \Ob(\Sch_{fppf})$ be a scheme. Let $F$ be an algebraic space over
$S$. A {\it small Zariski site $F_{Zar}$} of an algebraic space $F$ is one
of the sites described above.
\end{definition}

\noindent
Hence this gives a notion of what it means for something to be true
Zariski locally on an algebraic space, which is how we will use this
notion. In general the Zariski topology is not fine enough for our
purposes. For example we can consider the category of Zariski sheaves
on an algebraic space. It will turn out that this is not the
correct thing to consider, even for quasi-coherent sheaves.
One only gets the desired result when using the \'etale or fppf site of
$F$ to define quasi-coherent sheaves.











\section{Separation conditions on algebraic spaces}
\label{section-separation}

\noindent
A separation condition on an algebraic space $F$ is a condition
on the diagonal morphism $F \to F \times F$. Let us first
list the properties the diagonal has automatically.
Since the diagonal is representable by definition the following lemma
makes sense (through
Definition \ref{definition-relative-representable-property}).

\begin{lemma}
\label{lemma-properties-diagonal}
Let $S$ be a scheme contained in $\Sch_{fppf}$.
Let $F$ be an algebraic space over $S$.
Let $\Delta : F \to F \times F$ be the diagonal morphism.
Then
\begin{enumerate}
\item $\Delta$ is locally of finite type,
\item $\Delta$ is a monomorphism,
\item $\Delta$ is separated, and
\item $\Delta$ is locally quasi-finite.
\end{enumerate}
\end{lemma}

\begin{proof}
Let $F = U/R$ be a presentation of $F$.
As in the proof of Lemma \ref{lemma-presentation-quasi-compact} the diagram
$$
\xymatrix{
R \ar[r] \ar[d]_j & F \ar[d]^\Delta \\
U \times_S U \ar[r] & F \times F
}
$$
is cartesian. Hence according to
Lemma \ref{lemma-representable-morphisms-spaces-property}
it suffices to show that $j$ has the properties listed in the lemma.
(Note that each of the properties (1) -- (4) occur in the lists
of Remarks \ref{remark-list-properties-stable-base-change}
and \ref{remark-list-properties-fpqc-local-base}.)
Since $j$ is an equivalence relation it is a monomorphism.
Hence it is separated by
Schemes, Lemma \ref{schemes-lemma-monomorphism-separated}.
As $R$ is an \'etale equivalence relation we see that
$s, t : R \to U$ are \'etale. Hence $s, t$ are locally of finite
type. Then it follows from
Morphisms, Lemma \ref{morphisms-lemma-permanence-finite-type} that
$j$ is locally of finite type. Finally, as it is a monomorphism
its fibres are finite. Thus we conclude that it is locally quasi-finite by
Morphisms, Lemma \ref{morphisms-lemma-finite-fibre}.
\end{proof}

\noindent
Here are some common types of separation conditions, relative to the base
scheme $S$. There is also an absolute notion of these conditions which we
will discuss in
Properties of Spaces, Section \ref{spaces-properties-section-separation}.
Moreover, we will discuss separation conditions for a morphism of
algebraic spaces in
Morphisms of Spaces, Section \ref{spaces-morphisms-section-separation-axioms}.

\begin{definition}
\label{definition-separated}
Let $S$ be a scheme contained in $\Sch_{fppf}$.
Let $F$ be an algebraic space over $S$.
Let $\Delta : F \to F \times F$ be the diagonal morphism.
\begin{enumerate}
\item We say $F$ is {\it separated over $S$} if $\Delta$ is a closed immersion.
\item We say $F$ is {\it locally separated over $S$}\footnote{In the
literature this often refers to quasi-separated and
locally separated algebraic spaces.} if $\Delta$ is an
immersion.
\item We say $F$ is {\it quasi-separated over $S$} if $\Delta$ is quasi-compact.
\item We say $F$ is {\it Zariski locally quasi-separated over $S$}\footnote{This
definition was suggested by B.\ Conrad.} if there
exists a Zariski covering $F = \bigcup_{i \in I} F_i$ such that
each $F_i$ is quasi-separated.
\end{enumerate}
\end{definition}

\noindent
Note that if the diagonal is quasi-compact (when $F$ is separated or
quasi-separated) then the diagonal is actually
quasi-finite and separated, hence quasi-affine (by More on Morphisms,
Lemma \ref{more-morphisms-lemma-quasi-finite-separated-quasi-affine}).








\section{Examples of algebraic spaces}
\label{section-examples}

\noindent
In this section we construct some examples of algebraic spaces.
Some of these were suggested by B.\ Conrad.
Since we do not yet have a lot of theory at our disposal the
discussion is a bit awkward in some places.

\begin{example}
\label{example-affine-line-involution}
Let $k$ be a field of characteristic $\not = 2$. Let $U = \mathbf{A}^1_k$. Set
$$
j : R = \Delta \amalg \Gamma \longrightarrow U \times_k U
$$
where $\Delta = \{(x, x) \mid x \in \mathbf{A}^1_k\}$ and
$\Gamma = \{(x, -x) \mid x \in \mathbf{A}^1_k, x \not = 0\}$.
It is clear that $s, t : R \to U$ are \'etale, and hence
$j$ is an \'etale equivalence relation.
The quotient $X = U/R$ is an algebraic space by
Theorem \ref{theorem-presentation}.
Since $R$ is quasi-compact we see that $X$ is quasi-separated.
On the other hand, $X$ is not locally separated because
the morphism $j$ is not an immersion.
\end{example}

\begin{example}
\label{example-non-representable-descent}
Let $k$ be a field. Let $k'/k$ be a degree $2$ Galois extension
with $\text{Gal}(k'/k) = \{1, \sigma\}$. Let $S = \Spec(k[x])$
and $U = \Spec(k'[x])$. Note that
$$
U \times_S U =
\Spec((k' \otimes_k k')[x]) =
\Delta(U) \amalg \Delta'(U)
$$
where $\Delta' = (1, \sigma) : U \to U \times_S U$. Take
$$
R = \Delta(U) \amalg \Delta'(U \setminus \{0_U\})
$$
where $0_U \in U$ denotes the $k'$-rational point whose $x$-coordinate is zero.
It is easy to see that $R$ is an \'etale equivalence relation on $U$ over $S$
and hence $X = U/R$ is an algebraic space by
Theorem \ref{theorem-presentation}. Here are some properties of $X$ (some
of which will not make sense until later):
\begin{enumerate}
\item $X \to S$ is an isomorphism over $S \setminus \{0_S\}$,
\item the morphism $X \to S$ is \'etale (see
Properties of Spaces,
Definition \ref{spaces-properties-definition-etale})
\item the fibre $0_X$ of $X \to S$ over $0_S$ is isomorphic to
$\Spec(k') = 0_U$,
\item $X$ is not a scheme because if it were, then $\mathcal{O}_{X, 0_X}$
would be a local domain $(\mathcal{O}, \mathfrak m, \kappa)$ with
fraction field $k(x)$, with $x \in \mathfrak m$ and residue field
$\kappa = k'$ which is impossible,
\item $X$ is not separated, but it is
locally separated and quasi-separated,
\item there exists a surjective, finite, \'etale morphism $S' \to S$
such that the base change $X' = S' \times_S X$ is a scheme (namely, if
we base change to $S' = \Spec(k'[x])$ then $U$ splits into
two copies of $S'$ and $X'$ becomes isomorphic to the affine line with
$0$ doubled, see
Schemes, Example \ref{schemes-example-affine-space-zero-doubled}), and
\item if we think of $X$ as a finite type algebraic space over
$\Spec(k)$, then similarly the base change $X_{k'}$ is a scheme
but $X$ is not a scheme.
\end{enumerate}
In particular, this gives an example of a descent datum for schemes
relative to the covering $\{\Spec(k') \to \Spec(k)\}$
which is not effective.
\end{example}

\noindent
See also Examples, 
Lemma \ref{examples-lemma-non-effective-descent-projective},
which shows that descent data need not be effective
even for a projective morphism of schemes. That example
gives a smooth separated algebraic space of dimension 3
over ${\mathbf C}$ which is not a scheme.

\medskip\noindent
We will use the following lemma as a convenient way to construct
algebraic spaces as quotients of schemes by free group actions.

\begin{lemma}
\label{lemma-quotient}
Let $U \to S$ be a morphism of $\Sch_{fppf}$.
Let $G$ be an abstract group. Let $G \to \text{Aut}_S(U)$
be a group homomorphism. Assume
\begin{itemize}
\item[(*)] if $u \in U$ is a point, and $g(u) = u$
for some non-identity element $g \in G$, then $g$
induces a nontrivial automorphism of $\kappa(u)$.
\end{itemize}
Then
$$
j :
R = \coprod\nolimits_{g \in G} U
\longrightarrow
U \times_S U,
\quad
(g, x) \longmapsto (g(x), x)
$$
is an \'etale equivalence relation and hence
$$
F = U/R
$$
is an algebraic space by Theorem \ref{theorem-presentation}.
\end{lemma}

\begin{proof}
In the statement of the lemma the symbol $\text{Aut}_S(U)$ denotes
the group of automorphisms of $U$ over $S$.
Assume $(*)$ holds. Let us show that
$$
j :
R = \coprod\nolimits_{g \in G} U
\longrightarrow
U \times_S U,
\quad
(g, x) \longmapsto (g(x), x)
$$
is a monomorphism. This signifies that if $T$ is a nonempty
scheme, and $h : T \to U$ is a $T$-valued point such that
$g \circ h = g' \circ h$ then $g = g'$. Suppose
$T \not = \emptyset$, $h : T \to U$ and $g \circ h = g' \circ h$.
Let $t \in T$. Consider the composition
$\Spec(\kappa(t)) \to \Spec(\kappa(h(t))) \to U$.
Then we conclude that $g^{-1} \circ g'$ fixes $u = h(t)$ and
acts as the identity on its residue field. Hence $g = g'$ by $(*)$.

\medskip\noindent
Thus if $(*)$ holds we see that $j$ is a relation (see
Groupoids, Definition \ref{groupoids-definition-equivalence-relation}).
Moreover, it is an equivalence relation since on $T$-valued points
for a connected scheme $T$ we see that
$R(T) = G \times U(T) \to U(T) \times U(T)$ (recall that we always
work over $S$). Moreover, the morphisms $s, t : R \to U$ are \'etale
since $R$ is a disjoint product of copies of $U$.
This proves that $j : R \to U \times_S U$ is an \'etale equivalence relation.
\end{proof}

\noindent
Given a scheme $U$ and an action of a group $G$ on $U$ we say the action
of $G$ on $U$ is {\it free} if condition $(*)$ of Lemma \ref{lemma-quotient}
holds. This is equivalent to the notion of a free action of the constant
group scheme $G_S$ on $U$ as defined in
Groupoids, Definition \ref{groupoids-definition-free-action}.
The lemma can be interpreted as saying that quotients of schemes by
free actions of groups exist in the category of algebraic spaces.

\begin{definition}
\label{definition-quotient}
Notation $U \to S$, $G$, $R$ as in Lemma \ref{lemma-quotient}.
If the action of $G$ on $U$ satisfies $(*)$ we say $G$ {\it acts freely}
on the scheme $U$. In this case the algebraic space $U/R$ is denoted
$U/G$ and is called the {\it quotient of $U$ by $G$}.
\end{definition}

\noindent
This notation is consistent with the notation $U/G$ introduced in
Groupoids, Definition \ref{groupoids-definition-quotient-sheaf}.
We will later make sense of the quotient as an algebraic stack without
any assumptions on the action whatsoever; when we do this we will use the
notation $[U/G]$. Before we discuss the examples we prove
some more lemmas to facilitate the discussion. Here is a lemma discussing the
various separation conditions for this quotient when $G$ is finite.

\begin{lemma}
\label{lemma-quotient-finite-separated}
Notation and assumptions as in Lemma \ref{lemma-quotient}.
Assume $G$ is finite. Then
\begin{enumerate}
\item if $U \to S$ is quasi-separated, then $U/G$ is quasi-separated
over $S$, and
\item if $U \to S$ is separated, then $U/G$ is separated over $S$.
\end{enumerate}
\end{lemma}

\begin{proof}
In the proof of Lemma \ref{lemma-properties-diagonal}
we saw that it suffices to prove the
corresponding properties for the morphism $j : R \to U \times_S U$.
If $U \to S$ is quasi-separated, then for every affine open $V \subset U$
which maps into an affine of $S$
the opens $g(V) \cap V$ are quasi-compact. It follows that $j$ is
quasi-compact.
If $U \to S$ is separated, the diagonal $\Delta_{U/S}$ is a closed
immersion. Hence $j : R \to U \times_S U$ is a finite coproduct
of closed immersions with disjoint images. Hence $j$ is a closed immersion.
\end{proof}

\begin{lemma}
\label{lemma-quotient-field-map}
Notation and assumptions as in Lemma \ref{lemma-quotient}.
If $\Spec(k) \to U/G$ is a morphism, then there exist
\begin{enumerate}
\item a finite Galois extension $k'/k$,
\item a finite subgroup $H \subset G$,
\item an isomorphism $H \to \text{Gal}(k'/k)$, and
\item an $H$-equivariant morphism $\Spec(k') \to U$.
\end{enumerate}
Conversely, such data determine a morphism $\Spec(k) \to U/G$.
\end{lemma}

\begin{proof}
Consider the fibre product $V = \Spec(k) \times_{U/G} U$.
Here is a diagram
$$
\xymatrix{
V \ar[r] \ar[d] & U \ar[d] \\
\Spec(k) \ar[r] & U/G
}
$$
Then $V$ is a nonempty scheme \'etale over $\Spec(k)$ and hence is a
disjoint union $V = \coprod_{i \in I} \Spec(k_i)$
of spectra of fields $k_i$ finite separable over $k$
(Morphisms, Lemma \ref{morphisms-lemma-etale-over-field}).
We have
\begin{align*}
V \times_{\Spec(k)} V
& =
(\Spec(k) \times_{U/G} U) \times_{\Spec(k)}(\Spec(k) \times_{U/G} U) \\
& = 
\Spec(k) \times_{U/G} U \times_{U/G} U \\
& =
\Spec(k) \times_{U/G} U \times G \\
& =
V \times G
\end{align*}
The action of $G$ on $U$ induces an action of $a : G \times V \to V$.
The displayed equality means that
$G \times V \to V \times_{\Spec(k)} V$, $(g, v) \mapsto (a(g, v), v)$
is an isomorphism. In particular we see that for every $i$ we have
an isomorphism $H_i \times \Spec(k_i) \to \Spec(k_i \otimes_k k_i)$
where $H_i \subset G$ is the subgroup of elements fixing $i \in I$.
Thus $H_i$ is finite and is the Galois group of $k_i/k$.
We omit the converse construction.
\end{proof}

\noindent
It follows from this lemma for example that
if $k'/k$ is a finite Galois extension, then
$\Spec(k')/\text{Gal}(k'/k) \cong \Spec(k)$.
What happens if the extension is infinite? Here is an example.

\begin{example}
\label{example-Qbar}
Let $S = \Spec(\mathbf{Q})$.
Let $U = \Spec(\overline{\mathbf{Q}})$.
Let $G = \text{Gal}(\overline{\mathbf{Q}}/\mathbf{Q})$ with obvious
action on $U$. Then by construction property $(*)$ of
Lemma \ref{lemma-quotient} holds and we obtain an algebraic space
$$
X = \Spec(\overline{\mathbf{Q}})/G
\longrightarrow
S = \Spec(\mathbf{Q}).
$$
Of course this is totally ridiculous as an approximation of $S$!
Namely, by the Artin-Schreier theorem,
see \cite[Theorem 17, page 316]{JacobsonIII},
the only finite subgroups of $\text{Gal}(\overline{\mathbf{Q}}/\mathbf{Q})$
are $\{1\}$ and the conjugates of the order two group
$\text{Gal}(\overline{\mathbf{Q}}/\overline{\mathbf{Q}} \cap \mathbf{R})$.
Hence, if
$\Spec(k) \to X$ is a morphism with $k$ algebraic over $\mathbf{Q}$,
then it follows from Lemma \ref{lemma-quotient-field-map} and the theorem
just mentioned that either $k$ is $\overline{\mathbf{Q}}$ or isomorphic to
$\overline{\mathbf{Q}} \cap \mathbf{R}$.
\end{example}

\noindent
What is wrong with the example above is that
the Galois group comes equipped with a topology,
and this should somehow be part of any construction
of a quotient of $\Spec(\overline{\mathbf{Q}})$.
The following example is much more reasonable in my opinion
and may actually occur in ``nature''.

\begin{example}
\label{example-affine-line-translation}
Let $k$ be a field of characteristic zero.
Let $U = \mathbf{A}^1_k$ and let $G = \mathbf{Z}$.
As action we take $n(x) = x + n$, i.e., the action of
$\mathbf{Z}$ on the affine line by translation.
The only fixed point is the generic point and it
is clearly the case that $\mathbf{Z}$ injects into
the automorphism group of the field $k(x)$. (This is
where we use the characteristic zero assumption.)
Consider the morphism
$$
\gamma : \Spec(k(x)) \longrightarrow X = \mathbf{A}^1_k/\mathbf{Z}
$$
of the generic point of the affine line into the quotient.
We claim that this morphism does not factor through any
monomorphism $\Spec(L) \to X$ of the spectrum of
a field to $X$. (Contrary to what happens for schemes, see
Schemes, Section \ref{schemes-section-points}.) In fact, since
$\mathbf{Z}$ does not have any nontrivial finite subgroups we see from
Lemma \ref{lemma-quotient-field-map} that for any such
factorization $k(x) = L$. Finally, $\gamma$ is not a monomorphism
since
$$
\Spec(k(x)) \times_{\gamma, X, \gamma} \Spec(k(x))
\cong
\Spec(k(x)) \times \mathbf{Z}.
$$
\end{example}

\noindent
This example suggests that in order to define points of an algebraic space
$X$ we should consider equivalence classes of morphisms from spectra
of fields into $X$ and not the set of monomorphisms from spectra of fields.

\medskip\noindent
We finish with a truly awful example.

\begin{example}
\label{example-infinite-product}
Let $k$ be a field.
Let $A = \prod_{n \in \mathbf{N}} k$ be the infinite product.
Set $U = \Spec(A)$ seen as a scheme over $S = \Spec(k)$.
Note that the projection maps $\text{pr}_n : A \to k$ define open
and closed immersions $f_n : S \to U$. Set
$$
R = U \amalg \coprod\nolimits_{(n, m) \in \mathbf{N}^2, \ n \not = m} S
$$
with morphism $j$ equal to $\Delta_{U/S}$ on the component $U$
and $j = (f_n, f_m)$ on the component $S$ corresponding to $(n, m)$.
It is clear from the remark above that $s, t$ are \'etale.
It is also clear that $j$ is an equivalence relation. Hence we
obtain an algebraic space
$$
X = U/R.
$$
To see what this means we specialize to the case where
the field $k$ is finite with $q$ elements. Let us first
discuss the topological space $|U|$ associated to the scheme $U$
a little bit. All elements of $A$ satisfy $x^q = x$.
Hence every residue field of $A$ is isomorphic to $k$, and
all points of $U$ are closed. But the topology on $U$ isn't
the discrete topology. Let $u_n \in |U|$ be the point corresponding
to $f_n$. As mentioned above the points $u_n$ are
the open points (and hence isolated). This implies there have
to be other points since we know $U$ is quasi-compact, see
Algebra, Lemma \ref{algebra-lemma-quasi-compact}
(hence not equal to an infinite discrete set).
Another way to see this is because the (proper) ideal
$$
I =
\{x = (x_n) \in A \mid \text{all but a finite number of }x_n\text{ are zero}\}
$$
is contained in a maximal ideal. Note also that every element
$x$ of $A$ is of the form $x = ue$ where $u$ is a unit and $e$ is an
idempotent. Hence a basis for the topology of $A$ consists of open and
closed subsets (see Algebra, Lemma \ref{algebra-lemma-idempotent-spec}.)
So the topology on $|U|$ is totally disconnected, but nontrivial.
Finally, note that $\{u_n\}$ is dense in $|U|$.

\medskip\noindent
We will later define a topological space $|X|$ associated to $X$, see
Properties of Spaces, Section \ref{spaces-properties-section-points}.
What can we say about $|X|$?
It turns out that the map $|U| \to |X|$ is surjective and continuous.
All the points $u_n$ map to the same point $x_0$ of $|X|$, and none of
the other points get identified. Since $\{u_n\}$ is dense in $|U|$ we
conclude that the closure of $x_0$ in $|X|$ is $|X|$. In other words
$|X|$ is irreducible and $x_0$ is a generic point of $|X|$. This seems
bizarre since also $x_0$ is the image of a section
$S \to X$ of the structure morphism $X \to S$ (and in the case of
schemes this would imply it was a closed point, see
Morphisms, Lemma
\ref{morphisms-lemma-algebraic-residue-field-extension-closed-point-fibre}).
\end{example}

\noindent
Whatever you think is actually going on in this example, it certainly
shows that some care has to be exercised when defining irreducible
components, connectedness, etc of algebraic spaces.






\section{Change of big site}
\label{section-change-big-site}

\noindent
In this section we briefly discuss what happens when we change big sites.
The upshot is that we can always enlarge the big site at will, hence we
may assume any set of schemes we want to consider is contained in the big
fppf site over which we consider our algebraic space. Here is a precise
statement of the result.

\begin{lemma}
\label{lemma-change-big-site}
Suppose given big sites $\Sch_{fppf}$ and $\Sch'_{fppf}$.
Assume that $\Sch_{fppf}$ is contained in $\Sch'_{fppf}$,
see Topologies, Section \ref{topologies-section-change-alpha}.
Let $S$ be an object of $\Sch_{fppf}$. Let
\begin{align*}
g : \Sh((\Sch/S)_{fppf})
\longrightarrow
\Sh((\Sch'/S)_{fppf}), \\
f : \Sh((\Sch'/S)_{fppf})
\longrightarrow
\Sh((\Sch/S)_{fppf})
\end{align*}
be the morphisms of topoi of
Topologies, Lemma \ref{topologies-lemma-change-alpha}.
Let $F$ be a sheaf of sets on $(\Sch/S)_{fppf}$. Then
\begin{enumerate}
\item if $F$ is representable by a scheme
$X \in \Ob((\Sch/S)_{fppf})$ over $S$,
then $f^{-1}F$ is representable too, in fact it is representable by the
same scheme $X$, now viewed as an object of $(\Sch'/S)_{fppf}$, and
\item if $F$ is an algebraic space over $S$, then $f^{-1}F$ is an algebraic
space over $S$ also.
\end{enumerate}
\end{lemma}

\begin{proof}
Let $X \in \Ob((\Sch/S)_{fppf})$. Let us write $h_X$ for the
representable sheaf on $(\Sch/S)_{fppf}$ associated to $X$, and
$h'_X$ for the representable sheaf on $(\Sch'/S)_{fppf}$ associated to
$X$. By the description of $f^{-1}$ in
Topologies, Section \ref{topologies-section-change-alpha}
we see that $f^{-1}h_X = h'_X$. This proves (1).

\medskip\noindent
Next, suppose that $F$
is an algebraic space over $S$. By Lemma \ref{lemma-space-presentation}
this means that $F = h_U/h_R$ for some \'etale equivalence relation
$R \to U \times_S U$ in $(\Sch/S)_{fppf}$. Since $f^{-1}$ is an
exact functor we conclude that $f^{-1}F = h'_U/h'_R$. Hence
$f^{-1}F$ is an algebraic space over $S$ by Theorem \ref{theorem-presentation}.
\end{proof}

\noindent
Note that this lemma is purely set theoretical and has virtually no content.
Moreover, it is not true (in general) that the restriction of an algebraic
space over the bigger site is an algebraic space over the smaller site (simply
by reasons of cardinality). Hence we can only ever use a simple lemma of this
kind to enlarge the base category and never to shrink it.

\begin{lemma}
\label{lemma-fully-faithful}
Suppose $\Sch_{fppf}$ is contained in $\Sch'_{fppf}$.
Let $S$ be an object of $\Sch_{fppf}$. Denote
$\textit{Spaces}/S$ the category of algebraic spaces over $S$
defined using $\Sch_{fppf}$. Similarly, denote
$\textit{Spaces}'/S$ the category of algebraic spaces over $S$
defined using $\Sch'_{fppf}$. The construction of
Lemma \ref{lemma-change-big-site}
defines a fully faithful functor
$$
\textit{Spaces}/S \longrightarrow \textit{Spaces}'/S
$$
whose essential image consists of those $X' \in \Ob(\textit{Spaces}'/S)$
such that there exist $U, R \in \Ob((\Sch/S)_{fppf})$\footnote{Requiring the
existence of $R$ is necessary because of our choice of the function $Bound$ in
Sets, Equation (\ref{sets-equation-bound}). The size of the fibre product
$U \times_{X'} U$ can grow faster than $Bound$ in terms of the size of $U$. We
can illustrate this by setting $S = \Spec(A)$, $U = \Spec(A[x_i, i \in I])$ and
$R = \coprod_{(\lambda_i) \in A^I} \Spec(A[x_i, y_i]/(x_i - \lambda_i y_i))$.
In this case the size of $R$ grows like $\kappa^\kappa$ where $\kappa$ is the
size of $U$.} and morphisms
$$
U \longrightarrow X'
\quad\text{and}\quad
R \longrightarrow U \times_{X'} U
$$
in $\Sh((\Sch'/S)_{fppf})$ which are surjective as maps of sheaves
(for example if the displayed morphisms are surjective and \'etale).
\end{lemma}

\begin{proof}
In Sites, Lemma \ref{sites-lemma-bigger-site} we have seen that the functor
$f^{-1} : \Sh((\Sch/S)_{fppf}) \to \Sh((\Sch'/S)_{fppf})$
is fully faithful (see discussion in
Topologies, Section \ref{topologies-section-change-alpha}).
Hence we see that the displayed functor of the lemma is fully faithful.

\medskip\noindent
Suppose that $X' \in \Ob(\textit{Spaces}'/S)$ such that there exists
$U \in \Ob((\Sch/S)_{fppf})$ and a map $U \to X'$ in
$\Sh((\Sch'/S)_{fppf})$ which is surjective as a map of sheaves.
Let $U' \to X'$ be a surjective \'etale morphism with
$U' \in \Ob((\Sch'/S)_{fppf})$. Let $\kappa = \text{size}(U)$, see
Sets, Section \ref{sets-section-categories-schemes}.
Then $U$ has an affine open covering $U = \bigcup_{i \in I} U_i$
with $|I| \leq \kappa$. Observe that $U' \times_{X'} U \to U$ is \'etale
and surjective. For each $i$ we can pick a quasi-compact
open $U'_i \subset U'$ such that $U'_i \times_{X'} U_i \to U_i$
is surjective (because the scheme $U' \times_{X'} U_i$ is the
union of the Zariski opens $W \times_{X'} U_i$ for $W \subset U'$
affine and because $U' \times_{X'} U_i \to U_i$ is \'etale hence open).
Then $\coprod_{i \in I} U'_i \to X$ is surjective \'etale
because of our assumption that $U \to X$ and hence $\coprod U_i \to X$
is a surjection of sheaves (details omitted).
Because $U'_i \times_{X'} U \to U'_i$ is a surjection of sheaves
and because $U'_i$ is quasi-compact,
we can find a quasi-compact open $W_i \subset U'_i \times_{X'} U$
such that $W_i \to U'_i$ is surjective as a map of sheaves
(details omitted). Then $W_i \to U$ is \'etale and we conclude
that $\text{size}(W_i) \leq \text{size}(U)$, see
Sets, Lemma \ref{sets-lemma-bound-finite-type}. By
Sets, Lemma \ref{sets-lemma-bound-by-covering} we conclude
that $\text{size}(U'_i) \leq \text{size}(U)$.
Hence $\coprod_{i \in I} U'_i$ is isomorphic to an object of
$(\Sch/S)_{fppf}$ by Sets, Lemma \ref{sets-lemma-bound-size}.

\medskip\noindent
Now let $X'$, $U \to X'$ and $R \to U \times_{X'} U$ be as in the
statement of the lemma. In the previous paragraph we have seen that
we can find $U' \in \Ob((\Sch/S)_{fppf})$ and a surjective \'etale morphism
$U' \to X'$ in $\Sh((\Sch'/S)_{fppf})$. Then
$U' \times_{X'} U \to U'$ is a surjection of sheaves, i.e., we can find an fppf
covering $\{U'_i \to U'\}$ such that $U'_i \to U'$ factors through
$U' \times_{X'} U \to U'$.
By Sets, Lemma \ref{sets-lemma-bound-fppf-covering}
we can find
$\tilde U \to U'$
which is surjective, flat, and locally of finite presentation,
with $\text{size}(\tilde U) \leq \text{size}(U')$,
such that $\tilde U \to U'$ factors through
$U' \times_{X'} U \to U'$. Then we consider
$$
\xymatrix{
U' \times_{X'} U' \ar[d] &
\tilde U \times_{X'} \tilde U  \ar[l] \ar[d] \ar[r] &
U \times_{X'} U \ar[d] \\
U' \times_S U' & \tilde U \times_S \tilde U \ar[l] \ar[r] & U \times_S U
}
$$
The squares are cartesian. We know the objects of the bottom row
are represented by objects of $(\Sch/S)_{fppf}$. By the result of the
argument of the previous paragraph, the same is true for $U \times_{X'} U$
(as we have the surjection of sheaves $R \to U \times_{X'} U$
by assumption). Since $(\Sch/S)_{fppf}$ is closed under fibre
products (by construction), we see that $\tilde U \times_{X'} \tilde U$
is represented by an object of $(\Sch/S)_{fppf}$. Finally, the
map $\tilde U \times_{X'} \tilde U \to U' \times_{X'} U'$ is
a surjection of fppf sheaves as $\tilde U \to U'$ is so.
Thus we can once more apply the result of the previous paragraph
to conclude that $R' = U' \times_{X'} U'$ is represented by an object
of $(\Sch/S)_{fppf}$. At this point
Lemma \ref{lemma-space-presentation} and
Theorem \ref{theorem-presentation} imply that $X = h_{U'}/h_{R'}$
is an object of $\textit{Spaces}/S$ such that $f^{-1}X \cong X'$
as desired.
\end{proof}



\section{Change of base scheme}
\label{section-change-base-scheme}

\noindent
In this section we briefly discuss what happens when we change base schemes.
The upshot is that given a morphism $S \to S'$ of base schemes, any algebraic
space over $S$ can be viewed as an algebraic space over $S'$. And, given an
algebraic space $F'$ over $S'$ there is a base change $F'_S$ which is
an algebraic space over $S$.
We explain only what happens in case $S \to S'$ is a morphism of the
big fppf site under consideration, if only $S$ or $S'$ is contained in the
big site, then one first enlarges the big site as in
Section \ref{section-change-big-site}.

\begin{lemma}
\label{lemma-change-base-scheme}
Suppose given a big site $\Sch_{fppf}$.
Let $g : S \to S'$ be morphism of $\Sch_{fppf}$.
Let $j : (\Sch/S)_{fppf} \to (\Sch/S')_{fppf}$ be
the corresponding localization functor.
Let $F$ be a sheaf of sets on $(\Sch/S)_{fppf}$.
Then
\begin{enumerate}
\item for a scheme $T'$ over $S'$ we have
$j_!F(T'/S') =
\coprod\nolimits_{\varphi : T' \to S} F(T' \xrightarrow{\varphi} S),$
\item if $F$ is representable by a scheme
$X \in \Ob((\Sch/S)_{fppf})$,
then $j_!F$ is representable by $j(X)$ which is
$X$ viewed as a scheme over $S'$, and
\item if $F$ is an algebraic space over $S$, then $j_!F$ is an algebraic
space over $S'$, and if $F = U/R$ is a presentation, then
$j_!F = j(U)/j(R)$ is a presentation.
\end{enumerate}
Let $F'$ be a sheaf of sets on $(\Sch/S')_{fppf}$. Then
\begin{enumerate}
\item[(4)] for a scheme $T$ over $S$ we have $j^{-1}F'(T/S) = F'(T/S')$,
\item[(5)] if $F'$ is representable by a scheme
$X' \in \Ob((\Sch/S')_{fppf})$, then
$j^{-1}F'$ is representable, namely by $X'_S = S \times_{S'} X'$, and
\item[(6)] if $F'$ is an algebraic space, then
$j^{-1}F'$ is an algebraic space, and if $F' = U'/R'$ is a presentation,
then $j^{-1}F' = U'_S/R'_S$ is a presentation.
\end{enumerate}
\end{lemma}

\begin{proof}
The functors $j_!$, $j_*$ and $j^{-1}$ are defined in
Sites, Lemma \ref{sites-lemma-relocalize}
where it is also shown that $j = j_{S/S'}$ is the localization
of $(\Sch/S')_{fppf}$ at the object $S/S'$. Hence
all of the material on localization functors is available for $j$.
The formula in (1) is
Sites, Lemma \ref{sites-lemma-describe-j-shriek-good-site}.
By definition $j_!$ is the left adjoint to restriction $j^{-1}$,
hence $j_!$ is right exact. By
Sites, Lemma \ref{sites-lemma-j-shriek-commutes-equalizers-fibre-products}
it also commutes with fibre products and equalizers.
By
Sites, Lemma \ref{sites-lemma-describe-j-shriek-representable}
we see that $j_!h_X = h_{j(X)}$ hence (2) holds.
If $F$ is an algebraic space over $S$, then we can write $F = U/R$
(Lemma \ref{lemma-space-presentation})
and we get
$$
j_!F = j(U)/j(R)
$$
because $j_!$ being right exact commutes with coequalizers, and moreover
$j(R) = j(U) \times_{j_!F} j(U)$ as $j_!$ commutes with fibre products.
Since the morphisms $j(s), j(t) : j(R) \to j(U)$ are simply the morphisms
$s, t : R \to U$ (but viewed as morphisms of schemes over $S'$), they
are still \'etale. Thus $(j(U), j(R), s, t)$ is an \'etale equivalence relation.
Hence by
Theorem \ref{theorem-presentation}
we conclude that $j_!F$ is an algebraic space.

\medskip\noindent
Proof of (4), (5), and (6). The description of $j^{-1}$ is in
Sites, Section \ref{sites-section-localize}.
The restriction of the representable sheaf associated to $X'/S'$
is the representable sheaf associated to
$X'_S = S \times_{S'} Y'$ by
Sites, Lemma \ref{sites-lemma-localize-given-products}.
The restriction functor $j^{-1}$ is exact, hence $j^{-1}F' = U'_S/R'_S$.
Again by exactness the sheaf $R'_S$ is still an equivalence relation on
$U'_S$. Finally the two maps $R'_S \to U'_S$ are \'etale as base changes
of the \'etale morphisms $R' \to U'$. Hence $j^{-1}F' = U'_S/R'_S$ is
an algebraic space by
Theorem \ref{theorem-presentation}
and we win.
\end{proof}

\noindent
Note how the presentation $j_!F = j(U)/j(R)$ is just the presentation
of $F$ but viewed as a presentation by schemes over $S'$. Hence the
following definition makes sense.

\begin{definition}
\label{definition-base-change}
Let $\Sch_{fppf}$ be a big fppf site.
Let $S \to S'$ be a morphism of this site.
\begin{enumerate}
\item If $F'$ is an algebraic space over $S'$, then the
{\it base change of $F'$ to $S$} is the
algebraic space $j^{-1}F'$ described in
Lemma \ref{lemma-change-base-scheme}. We denote it $F'_S$.
\item If $F$ is an algebraic space over $S$, then $F$
{\it viewed as an algebraic space over $S'$}
is the algebraic space $j_!F$ over $S'$ described in
Lemma \ref{lemma-change-base-scheme}. We often simply denote this
$F$; if not then we will write $j_!F$.
\end{enumerate}
\end{definition}

\noindent
The algebraic space $j_!F$ comes equipped with a canonical morphism
$j_!F \to S$ of algebraic spaces over $S'$. This is true simply
because the sheaf $j_!F$ maps to $h_S$ (see for example the
explicit description in
Lemma \ref{lemma-change-base-scheme}).
In fact, in
Sites, Lemma \ref{sites-lemma-essential-image-j-shriek}
we have seen that the category of sheaves on $(\Sch/S)_{fppf}$
is equivalent to the category of pairs $(\mathcal{F}', \mathcal{F}' \to h_S)$
consisting of a sheaf on $(\Sch/S')_{fppf}$ and a map of sheaves
$\mathcal{F}' \to h_S$. The equivalence assigns to the sheaf $\mathcal{F}$
the pair $(j_!\mathcal{F}, j_!\mathcal{F} \to h_S)$.
This, combined with the above, leads to the following
result for categories of algebraic spaces.

\begin{lemma}
\label{lemma-category-of-spaces-over-smaller-base-scheme}
Let $\Sch_{fppf}$ be a big fppf site.
Let $S \to S'$ be a morphism of this site.
The construction above give an equivalence of
categories
$$
\left\{
\begin{matrix}
\text{category of algebraic}\\
\text{spaces over }S
\end{matrix}
\right\}
\leftrightarrow
\left\{
\begin{matrix}
\text{category of pairs }(F', F' \to S)\text{ consisting}\\
\text{of an algebraic space }F'\text{ over }S'\text{ and a}\\
\text{morphism }F' \to S\text{ of algebraic spaces over }S'
\end{matrix}
\right\}
$$
\end{lemma}

\begin{proof}
Let $F$ be an algebraic space over $S$. The functor from left to right
assigns the pair $(j_!F, j_!F \to S)$ to $F$
which is an object of the right hand side by
Lemma \ref{lemma-change-base-scheme}.
Since this defines an equivalence of categories of sheaves by
Sites, Lemma \ref{sites-lemma-essential-image-j-shriek}
to finish the proof it suffices to show:
if $F$ is a sheaf and $j_!F$ is an algebraic space, then $F$
is an algebraic space. To do this, write
$j_!F = U'/R'$ as in
Lemma \ref{lemma-space-presentation}
with $U', R' \in \Ob((\Sch/S')_{fppf})$.
Then the compositions $U' \to j_!F \to S$ and $R' \to j_!F \to S$
are morphisms of schemes over $S'$. Denote $U, R$ the corresponding
objects of $(\Sch/S)_{fppf}$. The two morphisms
$R' \to U'$ are morphisms over $S$ and hence correspond to
morphisms $R \to U$. Since these are simply the same
morphisms (but viewed over $S$) we see that we get an \'etale
equivalence relation over $S$. As $j_!$ defines an equivalence of
categories of sheaves (see reference above) we see that
$F = U/R$ and by
Theorem \ref{theorem-presentation}
we see that $F$ is an algebraic space.
\end{proof}

\noindent
The following lemma is a slight rephrasing of the above.

\begin{lemma}
\label{lemma-rephrase}
Let $\Sch_{fppf}$ be a big fppf site.
Let $S \to S'$ be a morphism of this site.
Let $F'$ be a sheaf on $(\Sch/S')_{fppf}$.
The following are equivalent:
\begin{enumerate}
\item The restriction $F'|_{(\Sch/S)_{fppf}}$
is an algebraic space over $S$, and
\item the sheaf $h_S \times F'$ is an algebraic space over $S'$.
\end{enumerate}
\end{lemma}

\begin{proof}
The restriction and the product match under
the equivalence of categories of
Sites, Lemma \ref{sites-lemma-essential-image-j-shriek}
so that
Lemma \ref{lemma-category-of-spaces-over-smaller-base-scheme}
above gives the result.
\end{proof}

\noindent
We finish this section with a lemma on a compatibility.

\begin{lemma}
\label{lemma-viewed-as-properties}
Let $\Sch_{fppf}$ be a big fppf site.
Let $S \to S'$ be a morphism of this site.
Let $F$ be an algebraic space over $S$.
Let $T$ be a scheme over $S$ and let $f : T \to F$ be
a morphism over $S$.
Let $f' : T' \to F'$ be the morphism over $S'$ we get from
$f$ by applying the equivalence of categories described in
Lemma \ref{lemma-category-of-spaces-over-smaller-base-scheme}.
For any property $\mathcal{P}$ as in
Definition \ref{definition-relative-representable-property}
we have $\mathcal{P}(f') \Leftrightarrow \mathcal{P}(f)$.
\end{lemma}

\begin{proof}
Suppose that $U$ is a scheme over $S$, and $U \to F$ is a surjective \'etale
morphism. Denote $U'$ the scheme $U$ viewed as a scheme over $S'$. In
Lemma \ref{lemma-change-base-scheme}
we have seen that $U' \to F'$ is surjective \'etale. Since
$$
j(T \times_{f, F} U) = T' \times_{f', F'} U'
$$
the morphism of schemes $T \times_{f, F} U \to U$
is identified with the morphism of schemes
$T' \times_{f', F'} U' \to U'$.
It is the same morphism, just viewed over different base schemes.
Hence the lemma follows from
Lemma \ref{lemma-representable-morphisms-spaces-property}.
\end{proof}










\begin{multicols}{2}[\section{Autres chapitres}]
\noindent
Préliminaires
\begin{enumerate}
\item \hyperref[introduction-section-phantom]{Introduction}
\item \hyperref[conventions-section-phantom]{Conventions}
\item \hyperref[sets-section-phantom]{Théorie des ensembles}
\item \hyperref[categories-section-phantom]{Catégories}
\item \hyperref[topology-section-phantom]{Topologie}
\item \hyperref[sheaves-section-phantom]{Faisceaux sur les espaces}
\item \hyperref[sites-section-phantom]{Sites et faisceaux}
\item \hyperref[stacks-section-phantom]{Champs}
\item \hyperref[fields-section-phantom]{Corps}
\item \hyperref[algebra-section-phantom]{Algèbre commutative}
\item \hyperref[brauer-section-phantom]{Groupes de Brauer}
\item \hyperref[homology-section-phantom]{Algèbre homologique}
\item \hyperref[derived-section-phantom]{Catégories dérivées}
\item \hyperref[simplicial-section-phantom]{Méthodes simpliciales}
\item \hyperref[more-algebra-section-phantom]{Compléments d'algèbre}
\item \hyperref[smoothing-section-phantom]{Lissage des morphismes d'anneaux}
\item \hyperref[modules-section-phantom]{Faisceaux de modules}
\item \hyperref[sites-modules-section-phantom]{Modules sur les sites}
\item \hyperref[injectives-section-phantom]{Objets injectifs}
\item \hyperref[cohomology-section-phantom]{Cohomologie des faisceaux}
\item \hyperref[sites-cohomology-section-phantom]{Cohomologie sur les sites}
\item \hyperref[dga-section-phantom]{Algèbre différentielle graduée}
\item \hyperref[dpa-section-phantom]{Algèbre à puissances divisées}
\item \hyperref[sdga-section-phantom]{Faisceaux différentiels gradués}
\item \hyperref[hypercovering-section-phantom]{Hyperrecouvrements}
\end{enumerate}
Schémas
\begin{enumerate}
\setcounter{enumi}{25}
\item \hyperref[schemes-section-phantom]{Schémas}
\item \hyperref[constructions-section-phantom]{Constructions de schémas}
\item \hyperref[properties-section-phantom]{Propriétés des schémas}
\item \hyperref[morphisms-section-phantom]{Morphismes de schémas}
\item \hyperref[coherent-section-phantom]{Cohomologie des schémas}
\item \hyperref[divisors-section-phantom]{Diviseurs}
\item \hyperref[limits-section-phantom]{Limites de schémas}
\item \hyperref[varieties-section-phantom]{Variétés}
\item \hyperref[topologies-section-phantom]{Topologies sur les schémas}
\item \hyperref[descent-section-phantom]{Descente}
\item \hyperref[perfect-section-phantom]{Catégories dérivées de schémas}
\item \hyperref[more-morphisms-section-phantom]{Compléments sur les morphismes}
\item \hyperref[flat-section-phantom]{Compléments sur la platitude}
\item \hyperref[groupoids-section-phantom]{Schémas en groupoïdes}
\item \hyperref[more-groupoids-section-phantom]{Compléments sur les schémas en groupoïdes}
\item \hyperref[etale-section-phantom]{Morphismes étales de schémas}
\end{enumerate}
Thèmes de théorie des schémas
\begin{enumerate}
\setcounter{enumi}{41}
\item \hyperref[chow-section-phantom]{Homologie de Chow}
\item \hyperref[intersection-section-phantom]{Théorie de l'intersection}
\item \hyperref[pic-section-phantom]{Schémas de Picard des courbes}
\item \hyperref[weil-section-phantom]{Théories cohomologiques de Weil}
\item \hyperref[adequate-section-phantom]{Modules adéquats}
\item \hyperref[dualizing-section-phantom]{Complexes dualisants}
\item \hyperref[duality-section-phantom]{Dualité pour les schémas}
\item \hyperref[discriminant-section-phantom]{Discriminants et différentes}
\item \hyperref[derham-section-phantom]{Cohomologie de de Rham}
\item \hyperref[local-cohomology-section-phantom]{Cohomologie locale}
\item \hyperref[algebraization-section-phantom]{Géométrie algébrique et formelle}
\item \hyperref[curves-section-phantom]{Courbes algébriques}
\item \hyperref[resolve-section-phantom]{Résolution des surfaces}
\item \hyperref[models-section-phantom]{Réduction semi-stable}
\item \hyperref[functors-section-phantom]{Foncteurs et morphismes}
\item \hyperref[equiv-section-phantom]{Catégories dérivées de variétés}
\item \hyperref[pione-section-phantom]{Groupes fondamentaux de schémas}
\item \hyperref[etale-cohomology-section-phantom]{Cohomologie étale}
\item \hyperref[crystalline-section-phantom]{Cohomologie cristalline}
\item \hyperref[proetale-section-phantom]{Cohomologie pro-étale}
\item \hyperref[relative-cycles-section-phantom]{Cycles relatifs}
\item \hyperref[more-etale-section-phantom]{Compléments de cohomologie étale}
\item \hyperref[trace-section-phantom]{La formule des traces}
\end{enumerate}
Espaces algébriques
\begin{enumerate}
\setcounter{enumi}{64}
\item \hyperref[spaces-section-phantom]{Espaces algébriques}
\item \hyperref[spaces-properties-section-phantom]{Propriétés des espaces algébriques}
\item \hyperref[spaces-morphisms-section-phantom]{Morphismes d'espaces algébriques}
\item \hyperref[decent-spaces-section-phantom]{Espaces algébriques décents}
\item \hyperref[spaces-cohomology-section-phantom]{Cohomologie des espaces algébriques}
\item \hyperref[spaces-limits-section-phantom]{Limites d'espaces algébriques}
\item \hyperref[spaces-divisors-section-phantom]{Diviseurs sur les espaces algébriques}
\item \hyperref[spaces-over-fields-section-phantom]{Espaces algébriques sur des corps}
\item \hyperref[spaces-topologies-section-phantom]{Topologies sur les espaces algébriques}
\item \hyperref[spaces-descent-section-phantom]{Descente et espaces algébriques}
\item \hyperref[spaces-perfect-section-phantom]{Catégories dérivées d'espaces}
\item \hyperref[spaces-more-morphisms-section-phantom]{Compléments sur les morphismes d'espaces}
\item \hyperref[spaces-flat-section-phantom]{Platitude sur les espaces algébriques}
\item \hyperref[spaces-groupoids-section-phantom]{Groupoïdes dans les espaces algébriques}
\item \hyperref[spaces-more-groupoids-section-phantom]{Compléments sur les groupoïdes dans les espaces}
\item \hyperref[bootstrap-section-phantom]{Amorçage}
\item \hyperref[spaces-pushouts-section-phantom]{Sommes amalgamées d'espaces algébriques}
\end{enumerate}
Thèmes de géométrie
\begin{enumerate}
\setcounter{enumi}{81}
\item \hyperref[spaces-chow-section-phantom]{Groupes de Chow des espaces}
\item \hyperref[groupoids-quotients-section-phantom]{Quotients de groupoïdes}
\item \hyperref[spaces-more-cohomology-section-phantom]{Compléments sur la cohomologie des espaces}
\item \hyperref[spaces-simplicial-section-phantom]{Espaces simpliciaux}
\item \hyperref[spaces-duality-section-phantom]{Dualité pour les espaces}
\item \hyperref[formal-spaces-section-phantom]{Espaces algébriques formels}
\item \hyperref[restricted-section-phantom]{Algébrisation des espaces formels}
\item \hyperref[spaces-resolve-section-phantom]{Résolution des surfaces revisitée}
\end{enumerate}
Théorie des déformations
\begin{enumerate}
\setcounter{enumi}{89}
\item \hyperref[formal-defos-section-phantom]{Théorie formelle des déformations}
\item \hyperref[defos-section-phantom]{Théorie des déformations}
\item \hyperref[cotangent-section-phantom]{Le complexe cotangent}
\item \hyperref[examples-defos-section-phantom]{Problèmes de déformation}
\end{enumerate}
Champs algébriques
\begin{enumerate}
\setcounter{enumi}{93}
\item \hyperref[algebraic-section-phantom]{Champs algébriques}
\item \hyperref[examples-stacks-section-phantom]{Exemples de champs}
\item \hyperref[stacks-sheaves-section-phantom]{Faisceaux sur les champs algébriques}
\item \hyperref[criteria-section-phantom]{Critères de représentabilité}
\item \hyperref[artin-section-phantom]{Axiomes d'Artin}
\item \hyperref[quot-section-phantom]{Espaces de Quot et de Hilbert}
\item \hyperref[stacks-properties-section-phantom]{Propriétés des champs algébriques}
\item \hyperref[stacks-morphisms-section-phantom]{Morphismes de champs algébriques}
\item \hyperref[stacks-limits-section-phantom]{Limites de champs algébriques}
\item \hyperref[stacks-cohomology-section-phantom]{Cohomologie des champs algébriques}
\item \hyperref[stacks-perfect-section-phantom]{Catégories dérivées de champs}
\item \hyperref[stacks-introduction-section-phantom]{Introduction aux champs algébriques}
\item \hyperref[stacks-more-morphisms-section-phantom]{Compléments sur les morphismes de champs}
\item \hyperref[stacks-geometry-section-phantom]{La géométrie des champs}
\end{enumerate}
Thèmes de théorie des espaces de modules
\begin{enumerate}
\setcounter{enumi}{107}
\item \hyperref[moduli-section-phantom]{Champs de modules}
\item \hyperref[moduli-curves-section-phantom]{Espaces de modules de courbes}
\end{enumerate}
Divers
\begin{enumerate}
\setcounter{enumi}{109}
\item \hyperref[examples-section-phantom]{Exemples}
\item \hyperref[exercises-section-phantom]{Exercices}
\item \hyperref[guide-section-phantom]{Guide bibliographique}
\item \hyperref[desirables-section-phantom]{Desiderata}
\item \hyperref[coding-section-phantom]{Style de programmation}
\item \hyperref[obsolete-section-phantom]{Obsolète}
\item \hyperref[fdl-section-phantom]{Licence de documentation libre GNU}
\item \hyperref[index-section-phantom]{Index généré automatiquement}
\end{enumerate}
\end{multicols}


\bibliography{my}
\bibliographystyle{amsalpha}

\end{document}
